% ============================================================
%  يُترجَم بـ XeLaTeX  (في Overleaf: Menu → Compiler → XeLaTeX)
%  الخط: Amiri (متوفر في Overleaf، أو حمّله من github.com/aliftype/amiri)
% ============================================================
\documentclass[12pt,a4paper]{article}
\usepackage[a4paper,margin=1.7cm,top=2.3cm,bottom=2.9cm,headheight=15pt,footskip=1.5cm]{geometry}
\usepackage{amsmath,amssymb}
\usepackage{icomma}
\usepackage[english,bidi=default,layout=lists]{babel}
\babelprovide[import,main]{arabic}
\babelfont{rm}[Scale=1.08]{Amiri}
\babelfont{sf}[Scale=1.08]{Amiri}
\usepackage[most]{tcolorbox}
\usepackage{tikz}
\usetikzlibrary{arrows.meta,calc,shadings,patterns}
\usepackage{array}
\usepackage{colortbl}
\usepackage{fancyhdr}
\usepackage{titlesec}
\usepackage{enumitem}
\usepackage{collcell}
\usepackage{needspace}
\usepackage{etoolbox}
\newcolumntype{C}{>{\collectcell\ar}c<{\endcollectcell}}
\pretocmd{\section}{\Needspace{15\baselineskip}}{}{}
\pretocmd{\subsection}{\Needspace{11\baselineskip}}{}{}

\definecolor{cDef}{HTML}{1F4E9C}
\definecolor{cThm}{HTML}{B3261E}
\definecolor{cRem}{HTML}{2E7D32}
\definecolor{cEx}{HTML}{C76B00}
\definecolor{cMeth}{HTML}{6A1B9A}
\definecolor{cSol}{HTML}{00695C}
\definecolor{cSec}{HTML}{1F3A68}

\titleformat{\section}{\Large\bfseries\color{cSec}}{\thesection.}{0.6em}{}[{\color{cSec!60}\titlerule[1pt]}]
\titleformat{\subsection}{\large\bfseries\color{cThm!85!black}}{\thesubsection.}{0.5em}{}
\titlespacing*{\section}{0pt}{14pt}{8pt}
\titlespacing*{\subsection}{0pt}{10pt}{5pt}
\setlength{\parindent}{0pt}
\setlength{\parskip}{4pt}
\setlist{itemsep=2pt,topsep=3pt}

\newtcolorbox{arbox}[3][]{enhanced,
  colback=#2!5,colframe=#2!80!black,arc=2mm,boxrule=0.8pt,
  left=6pt,right=6pt,top=8pt,bottom=5pt,
  fonttitle=\bfseries,coltitle=white,
  title={\foreignlanguage{arabic}{#3}},
  attach boxed title to top right={xshift=-5mm,yshift=-3mm},
  boxed title style={colback=#2!80!black,arc=1.5mm,boxrule=0pt},
  before upper={\selectlanguage{arabic}},
  before skip=8pt,after skip=8pt,breakable,#1}
\newenvironment{BOX}[2]{\begin{otherlanguage}{english}\begin{arbox}{#1}{#2}}{\end{arbox}\end{otherlanguage}}
\newenvironment{tarif}[1][تعريف]{\begin{BOX}{cDef}{#1}}{\end{BOX}}
\newenvironment{khassia}[1][خاصية]{\begin{BOX}{cThm}{#1}}{\end{BOX}}
\newenvironment{molahada}[1][ملاحظة]{\begin{BOX}{cRem}{#1}}{\end{BOX}}
\newenvironment{mithal}[1][مثال]{\begin{BOX}{cEx}{#1}}{\end{BOX}}
\newenvironment{tariqa}[1][طريقة]{\begin{BOX}{cMeth}{#1}}{\end{BOX}}

\newcounter{exo}
\newenvironment{tamrin}[1]{\stepcounter{exo}\begin{BOX}{cSec}{التمرين \theexo: #1}}{\end{BOX}}
\newenvironment{hal}{\Needspace{12\baselineskip}\begin{otherlanguage}{english}\begin{arbox}[breakable]{cSol}{حل التمرين \theexo}}{\end{arbox}\end{otherlanguage}}

\newenvironment{LTR}{\begin{otherlanguage}{english}}{\end{otherlanguage}}
\newcommand{\ar}[1]{\foreignlanguage{arabic}{#1}}

\renewcommand{\dots}{\mathinner{\ldotp\ldotp\ldotp}}
\renewcommand{\ldots}{\mathinner{\ldotp\ldotp\ldotp}}
\renewcommand{\cdots}{\mathinner{\cdotp\cdotp\cdotp}}
\newcommand{\N}{\mathbb{N}}
\newcommand{\Z}{\mathbb{Z}}
\newcommand{\D}{\mathbb{D}}
\newcommand{\Q}{\mathbb{Q}}
\newcommand{\R}{\mathbb{R}}
\newcommand{\pgcd}{\operatorname{PGCD}}
\newcommand{\ppcm}{\operatorname{PPCM}}

% \NLseg{a}{b}{[}{]} ; \NLright{a}{[} ; \NLleft{b}{]}  (interval on a number line)
\newcommand{\NLseg}[4]{\begin{LTR}\begin{tikzpicture}[baseline=-0.6ex]
 \draw[-{Stealth[length=2mm]},gray!70](0,0)--(3.4,0);
 \draw[line width=2pt,cThm!80!black](0.7,0)--(2.7,0);
 \node[cThm!80!black,font=\LARGE] at (0.7,0){$#3$};
 \node[cThm!80!black,font=\LARGE] at (2.7,0){$#4$};
 \node[below=3pt,font=\footnotesize] at (0.7,0){$#1$};
 \node[below=3pt,font=\footnotesize] at (2.7,0){$#2$};
\end{tikzpicture}\end{LTR}}
\newcommand{\NLright}[2]{\begin{LTR}\begin{tikzpicture}[baseline=-0.6ex]
 \draw[-{Stealth[length=2mm]},gray!70](0,0)--(3.4,0);
 \draw[line width=2pt,cThm!80!black](1.3,0)--(3.3,0);
 \node[cThm!80!black,font=\LARGE] at (1.3,0){$#2$};
 \node[below=3pt,font=\footnotesize] at (1.3,0){$#1$};
\end{tikzpicture}\end{LTR}}
\newcommand{\NLleft}[2]{\begin{LTR}\begin{tikzpicture}[baseline=-0.6ex]
 \draw[-{Stealth[length=2mm]},gray!70](0,0)--(3.4,0);
 \draw[line width=2pt,cThm!80!black](0.05,0)--(2.1,0);
 \node[cThm!80!black,font=\LARGE] at (2.1,0){$#2$};
 \node[below=3pt,font=\footnotesize] at (2.1,0){$#1$};
\end{tikzpicture}\end{LTR}}

% ---------- social-media icons (pure TikZ) ----------
\definecolor{cFB}{HTML}{1877F2}
\definecolor{cTG}{HTML}{229ED9}
\definecolor{cIG1}{HTML}{FEDA75}
\definecolor{cIG2}{HTML}{D62976}
\definecolor{cIG3}{HTML}{4F5BD5}
\newcommand{\fbicon}{\begin{LTR}\begin{tikzpicture}[baseline=-0.75ex,x=1em,y=1em]
 \fill[cFB] (0,0) circle (0.55);
 \node[white,inner sep=0,font=\usefont{TU}{lmss}{bx}{n}\small] at (0.04,-0.02) {f};
\end{tikzpicture}\end{LTR}}
\newcommand{\igicon}{\begin{LTR}\begin{tikzpicture}[baseline=-0.75ex,x=1em,y=1em]
 \shade[lower left=cIG1,lower right=cIG2,upper left=cIG2,upper right=cIG3,rounded corners=0.28em] (-0.55,-0.55) rectangle (0.55,0.55);
 \draw[white,line width=0.09em,rounded corners=0.18em] (-0.36,-0.36) rectangle (0.36,0.36);
 \draw[white,line width=0.09em] (0,0) circle (0.17);
 \fill[white] (0.22,0.22) circle (0.05);
\end{tikzpicture}\end{LTR}}
\newcommand{\tgicon}{\begin{LTR}\begin{tikzpicture}[baseline=-0.75ex,x=1em,y=1em]
 \fill[cTG] (0,0) circle (0.55);
 \fill[white] (-0.33,0.0) -- (0.32,0.27) -- (0.2,-0.3) -- (0.0,-0.15) -- (-0.08,-0.28) -- (-0.08,-0.08) -- cycle;
 \fill[cTG!70!black] (-0.08,-0.08) -- (0.25,0.2) -- (0.0,-0.15) -- (-0.08,-0.28) -- cycle;
\end{tikzpicture}\end{LTR}}
% page number: black box, white two-digit number, at the right of the social row
\newcommand{\PageBox}{\begin{LTR}\tikz[baseline=(p.base)]\node[fill=black,text=white,inner sep=2.5pt,minimum width=1.6em,minimum height=1.3em,font=\usefont{TU}{lmss}{bx}{n}\small](p){\ifnum\value{page}<10 0\fi\arabic{page}};\end{LTR}}
% visual left-to-right order: Telegram | Instagram | Facebook | page box
\newcommand{\SocialRow}{\begin{otherlanguage}{english}%
 \ar{\textbf{الأستاذ سوايسية هشام للرياضيات}}\ \tgicon\hfill
 \textbf{prof\_hicham\_math}\ \igicon\hfill
 \ar{\textbf{الأستاذ سوايسية هشام للرياضيات}}\ \fbicon\hspace{2.2em}\PageBox\end{otherlanguage}}

\pagestyle{fancy}
\fancyhf{}
\renewcommand{\headrulewidth}{0.6pt}
\renewcommand{\footrulewidth}{0.4pt}
% \SetHeader{right head: document title}{left head: level -- stream}
\newcommand{\SetHeader}[2]{%
 \fancyhead[C]{\parbox{\textwidth}{\ar{#1\hfill #2}}}%
 \fancyfoot[C]{\parbox{\textwidth}{\small\SocialRow}}}
\usepackage{tkz-tab}
\usepackage{multicol}
% variation / sign table (LTR)
\newenvironment{TAB}{\par\begin{center}\begin{LTR}\begin{tikzpicture}[scale=0.95,every node/.style={font=\small}]}{\end{tikzpicture}\end{LTR}\end{center}}
% grid + axes, clipped: \begin{GR}{xmin}{xmax}{ymin}{ymax}{unit cm}
\newenvironment{GR}[5]{\par\begin{center}\begin{LTR}\begin{tikzpicture}[x=#5cm,y=#5cm,font=\footnotesize]
 \draw[gray!25,very thin] (#1,#3) grid (#2,#4);
 \draw[-{Stealth[length=2mm]}] (#1,0)--(#2+0.3,0) node[below left]{$x$};
 \draw[-{Stealth[length=2mm]}] (0,#3)--(0,#4+0.3) node[below left]{$y$};
 \node[below left] at (0,0){$O$};
 \begin{scope}\clip (#1,#3) rectangle (#2,#4);}
{\end{scope}\end{tikzpicture}\end{LTR}\end{center}}
\newcommand{\Ticks}[2]{\foreach \t in {#1} \node[below,font=\scriptsize] at (\t,0){$\t$}; \foreach \t in {#2} \node[left,font=\scriptsize] at (0,\t){$\t$};}
\newcommand{\bx}[1]{\ensuremath{\colorbox{cDef!12}{$\displaystyle\color{cThm}#1$}}}
\newcommand{\LW}[1]{{\color{cDef}#1}}  % law / formula colour
\newcommand{\TwoCol}[4]{\par\begin{LTR}\noindent\begin{minipage}[t]{#1\textwidth}\selectlanguage{arabic}#2\end{minipage}\hfill\begin{minipage}[t]{#3\textwidth}\selectlanguage{arabic}#4\end{minipage}\end{LTR}\par}
\newcommand{\RL}[4]{\TwoCol{#3}{#4}{#1}{#2}}  % \RL{w1}{right content}{w2}{left content}
\newcommand{\ok}{$\textcolor{cRem}{\checkmark}$}
\newcommand{\limn}{\displaystyle\lim_{n\to+\infty}}
% sub-questions, each on its own line
\newenvironment{sq}{\begin{itemize}[leftmargin=2.4em,itemsep=3pt]}{\end{itemize}}
\newcommand{\qa}{\item[(أ)]}\newcommand{\qb}{\item[(ب)]}\newcommand{\qc}{\item[(جـ)]}\newcommand{\qd}{\item[(د)]}\newcommand{\qe}{\item[(هـ)]}
\newcommand{\pa}{\par\smallskip\textbf{(أ)} }\newcommand{\pb}{\par\smallskip\textbf{(ب)} }\newcommand{\pc}{\par\smallskip\textbf{(جـ)} }\newcommand{\pd}{\par\smallskip\textbf{(د)} }\newcommand{\pe}{\par\smallskip\textbf{(هـ)} }
\newcommand{\QQ}[1]{\par\smallskip\textbf{#1.} }
\newenvironment{tmr}{\stepcounter{exo}\begin{BOX}{cSec}{تمرين محلول \theexo}}{\end{BOX}}
\newcommand{\ABC}[3]{\\[2pt]\begin{tabular}{p{0.28\textwidth}p{0.28\textwidth}p{0.28\textwidth}}(جـ) \ #3 & (ب) \ #2 & (أ) \ #1\end{tabular}}
\newtcbox{\bxtbox}{on line,coltext=cThm,colback=cDef!12,colframe=cDef!12,boxrule=0pt,arc=1pt,boxsep=0pt,left=2pt,right=2pt,top=1pt,bottom=1pt}
\newcommand{\bxt}[1]{\foreignlanguage{english}{\bxtbox{\ar{#1}}}}
\SetHeader{\small\bfseries\color{cSec}تمارين الدالة اللوغاريتمية}{\small\color{cSec}السنة الثالثة ثانوي -- تسيير واقتصاد}
\newenvironment{FIG}{\par\begin{center}\begin{LTR}}{\end{LTR}\end{center}}
\renewcommand{\bx}[1]{\ensuremath{\colorbox{cDef!12}{$\displaystyle\color{cThm}#1$}}}
\newcommand{\Car}[1]{\left\{\begin{array}{l}#1\end{array}\right.}
\newcommand{\Lm}[1]{\displaystyle\lim_{x\to#1}}
\newcommand{\LD}[3]{\par لدينا: $\Lm{#2}#1=\bx{#3}$ لأنّ: \quad}
\newcommand{\E}{\par\smallskip}
\newcommand{\EXP}{\textcolor{cThm}{\textbf{(\ding{72}\ding{72})}}}
\usepackage{pifont}
\relpenalty=10000\binoppenalty=10000\tolerance=9000\emergencystretch=0.3\linewidth
\begin{document}
\begin{LTR}
\begin{tcolorbox}[enhanced,colback=cSec,colframe=cSec,arc=3mm,halign=center,boxrule=0pt,top=12pt,bottom=12pt,before upper={\selectlanguage{arabic}}]
{\color{white}\fontsize{28}{34}\selectfont\bfseries تمارين الدالة اللوغاريتمية}\\[8pt]
{\color{white}\large تمارين من الاختبارات + مسائل البكالوريا (2016 -- 2026) + اختيار من متعدد ، مع الحلول المفصلة}\\[5pt]
{\color{white}\large السنة الثالثة ثانوي -- شعبة تسيير واقتصاد}
\end{tcolorbox}
\end{LTR}

\begin{molahada}
\begin{itemize}
\item الأسئلة المشار إليها بـ $(\star)$ تتطلب درس الدالة الأصلية والحساب التكاملي ، ويمكن تأجيلها إلى ما بعد دراسته.
\item الأسئلة المشار إليها بـ \EXP\ تستعمل الدالة الأسية: في شعبة تسيير واقتصاد تُدرس الدالة اللوغاريتمية \textbf{قبل} الدالة الأسية ، فتؤجَّل إلى ما بعد درس الدالة الأسية.
\item النهايات مكتوبة على الشكل: «لدينا: النهاية $=$ النتيجة لأنّ:» ثم النهايات الجزئية في حاضنة.
\end{itemize}
\end{molahada}

% =====================================================
\section{تمارين من الاختبارات (البكالوريا التجريبية)}

\begin{tamrin}{بكالوريا تجريبية 2026 -- تسيير واقتصاد (بوقيرات) -- التمرين الثالث}
ليكن في $\R$ كثير الحدود $P(x)$ حيث: $P(x)=2x^3-x^2-4050x+2025$.
\QQ{1} تحقق أنه من أجل كل عدد حقيقي $x$: $P(x)=(2x-1)\left(x^2-2025\right)$ ، ثم حل في $\R$ المعادلة $P(x)=0$.
\QQ{2} استنتج حلول المعادلتين:
\begin{sq}
\qa $2(\log x)^3-(\log x)^2-4050\log x+2025=0$ حيث $x>0$.
\qb \EXP\ $2e^{3x}-e^{2x}-4050e^x+2025=0$.
\end{sq}
\QQ{3} حل في المجال $]0;+\infty[$ المتراجحة: $(2\ln x-1)\left((\ln x)^2-2025\right)\geqslant0$.
\end{tamrin}

\begin{hal}
\QQ{1} ننشر:
\[(2x-1)\left(x^2-2025\right)=2x^3-4050x-x^2+2025=\bx{P(x)}\]
$P(x)=0$ يعني $2x-1=0$ أو $x^2=2025$ ، أي $x=\dfrac12$ أو $x=45$ أو $x=-45$:
\[\bx{S=\left\{-45;\frac12;45\right\}}\]
\QQ{2}
\pa نضع $X=\log x$ ، فتصبح المعادلة $P(X)=0$ ، ومنه $\log x\in\left\{-45;\dfrac12;45\right\}$. وبما أن $\log x=y$ يعني $x=10^y$:
\[\bx{S=\left\{10^{-45};\sqrt{10};10^{45}\right\}}\]
\pb \EXP\ نضع $X=e^x>0$: $P(X)=0$ ، والحل $X=-45$ مرفوض ، ومنه $e^x=\dfrac12$ أو $e^x=45$:
\[\bx{S=\{-\ln2;\ln45\}}\]
\QQ{3} نضع $X=\ln x$: المتراجحة هي $P(X)\geqslant0$ لأن $(2X-1)(X^2-2025)=P(X)=(2X-1)(X-45)(X+45)$.
\begin{TAB}
\tkzTabInit[lgt=1.6,espcl=2.2]{$X$/1,$P(X)$/1}{$-\infty$,$-45$,$\frac12$,$45$,$+\infty$}
\tkzTabLine{,-,z,+,z,-,z,+,}
\end{TAB}
$P(X)\geqslant0$ يعني $-45\leqslant X\leqslant\dfrac12$ أو $X\geqslant45$ ، أي $-45\leqslant\ln x\leqslant\dfrac12$ أو $\ln x\geqslant45$:
\[\bx{S=\left[e^{-45};\sqrt e\right]\cup\left[e^{45};+\infty\right[}\]
\end{hal}

\begin{tamrin}{بكالوريا تجريبية 2026 -- تسيير واقتصاد (بوقيرات) -- التمرين الرابع}
\textbf{I)} الدالة $g$ معرفة على $]0;+\infty[$ بـ: $g(x)=3x^3-2+4\ln x$ ، وهي متزايدة تمامًا من $-\infty$ إلى $+\infty$.
\QQ{1} بيّن أن المعادلة $g(x)=0$ تقبل حلًا وحيدًا $\alpha$ حيث: $0{,}9<\alpha<1$.
\QQ{2} استنتج إشارة $g(x)$ حسب قيم $x$ من $]0;+\infty[$.\par
\textbf{II)} $f$ الدالة المعرفة على $]0;+\infty[$ بـ: $f(x)=3x-2-\dfrac{2\ln x}{x^2}$ ، و $(C_f)$ تمثيلها البياني (وحدة الطول $2\,cm$).
\QQ{1} أحسب $\Lm{0^+}f(x)$ و $\Lm{+\infty}f(x)$. (يعطى: $\Lm{+\infty}\dfrac{\ln x}{x^2}=0$).
\QQ{2}
\begin{sq}
\qa بيّن أن المستقيم $(\Delta)$ ذا المعادلة $y=3x-2$ مقارب مائل للمنحنى $(C_f)$.
\qb ادرس وضعية المنحنى $(C_f)$ بالنسبة إلى المستقيم $(\Delta)$.
\end{sq}
\QQ{3}
\begin{sq}
\qa بيّن أنه من أجل كل $x$ من $]0;+\infty[$: $f'(x)=\dfrac{g(x)}{x^3}$.
\qb استنتج اتجاه تغير الدالة $f$ ثم شكّل جدول تغيراتها.
\qc ليكن $(T)$ مماس $(C_f)$ عند النقطة ذات الفاصلة $x_0$ ، معادلته $y=3x-2-e^{-1}$. أحسب $x_0$.
\end{sq}
\QQ{4} ارسم كلًا من $(\Delta)$ و $(C_f)$. (نأخذ $f(\alpha)\approx0{,}9$).
\QQ{5} $(\star)$ $H$ الدالة المعرفة على $]0;+\infty[$ بـ: $H(x)=\dfrac{1+\ln x}{x}$.
\begin{sq}
\qa بيّن أن $H$ دالة أصلية للدالة $x\mapsto-\dfrac{\ln x}{x^2}$.
\qb أحسب بـ $cm^2$ مساحة الحيز المحدد بالمنحنى $(C_f)$ والمستقيم $(\Delta)$ والمستقيمين $x=1$ و $x=e$.
\end{sq}
\end{tamrin}

\begin{hal}
\textbf{I)}
\QQ{1} $g$ مستمرة ومتزايدة تمامًا على $[0{,}9;1]$ ، و:
\[g(0{,}9)=2{,}187-2+4\ln0{,}9\approx-0{,}23<0\qquad;\qquad g(1)=3-2+0=1>0\]
حسب مبرهنة القيم المتوسطة: المعادلة $g(x)=0$ تقبل حلًا وحيدًا $\bx{\alpha\in]0{,}9;1[}$ ($\alpha\approx0{,}92$).
\QQ{2} $g$ متزايدة تمامًا و $g(\alpha)=0$:
\begin{TAB}
\tkzTabInit[lgt=1.6,espcl=3]{$x$/1,$g(x)$/1}{$0$,$\alpha$,$+\infty$}
\tkzTabLine{d,-,z,+,}
\end{TAB}
\textbf{II)}
\QQ{1}
\LD{f(x)}{0^+}{+\infty} $\Car{\Lm{0^+}(3x-2)=-2\\[6pt]\Lm{0^+}\left(-\dfrac{2\ln x}{x^2}\right)=+\infty}$\E
لأن $\ln x\to-\infty$ و $x^2\to0^+$. إذن محور التراتيب مقارب عمودي.
\LD{f(x)}{+\infty}{+\infty} $\Car{\Lm{+\infty}(3x-2)=+\infty\\[6pt]\Lm{+\infty}\dfrac{2\ln x}{x^2}=0}$\E
\QQ{2}
\pa $f(x)-(3x-2)=-\dfrac{2\ln x}{x^2}$ و $\Lm{+\infty}\left(-\dfrac{2\ln x}{x^2}\right)=0$: \bxt{$(\Delta)$ مقارب مائل لـ $(C_f)$ بجوار $+\infty$}.
\pb إشارة الفرق عكس إشارة $\ln x$ ، وينعدم عند $x=1$ ($f(1)=1$):
\begin{FIG}
\begin{tikzpicture}[font=\normalsize]
\fill[gray!12] (0,0) rectangle (2.6,-0.9);
\fill[green!18] (2.6,0) rectangle (15,-0.9);
\fill[orange!45] (0,-0.9) rectangle (2.6,-4.6);
\fill[cyan!22] (2.6,-0.9) rectangle (15,-4.6);
\draw[thick] (0,0) rectangle (15,-4.6); \draw (0,-0.9)--(15,-0.9) (0,-1.8)--(15,-1.8) (2.6,0)--(2.6,-4.6);
\node at (1.3,-0.45){$x$};
\node at (1.3,-1.35){$-\frac{2\ln x}{x^2}$};
\node[align=center,font=\bfseries] at (1.3,-3.2){\ar{الوضع}\\ \ar{النسبي}};
\node at (3.3,-0.45){$0$};
\node at (8.8,-0.45){$1$};
\node at (14.3,-0.45){$+\infty$};
\node at (5.3,-1.35){$+$};
\node[align=center,font=\small] at (5.3,-3.2){$(C_f)$\\ \ar{فوق}\\ $(\Delta)$};
\node at (12.3,-1.35){$-$};
\node[align=center,font=\small] at (12.3,-3.2){$(C_f)$\\ \ar{تحت}\\ $(\Delta)$};
\draw[double,thick] (3.3,-0.9)--(3.3,-4.6);
\node at (8.8,-1.35){$0$};
\draw[densely dotted] (8.8,-0.9)--(8.8,-1.13) (8.8,-1.57)--(8.8,-1.8);
\draw (7.3,-4.6) .. controls (7.375,-3.06) and (8.35,-2.136) .. (8.8,-1.8) .. controls (9.25,-2.136) and (10.225,-3.06) .. (10.3,-4.6);
\node[align=center,font=\footnotesize,inner sep=0pt] at (8.8,-3.65){\ar{$(C_f)$ يقطع}\\ $(\Delta)$\\ \ar{في النقطة}\\ $(1;1)$};
\end{tikzpicture}
\end{FIG}
\QQ{3}
\pa $\left(\dfrac{\ln x}{x^2}\right)'=\dfrac{\frac1x\times x^2-2x\ln x}{x^4}=\dfrac{1-2\ln x}{x^3}$ ، ومنه:
\[f'(x)=3-\frac{2(1-2\ln x)}{x^3}=\frac{3x^3-2+4\ln x}{x^3}=\bx{\frac{g(x)}{x^3}}\]
\pb إشارة $f'(x)$ هي إشارة $g(x)$ (لأن $x^3>0$): \bxt{$f$ متناقصة تمامًا على $]0;\alpha]$ ومتزايدة تمامًا على $[\alpha;+\infty[$}.
\begin{TAB}
\tkzTabInit[lgt=1.6,espcl=3]{$x$/1,$f'(x)$/1,$f(x)$/2}{$0$,$\alpha$,$+\infty$}
\tkzTabLine{d,-,z,+,}
\tkzTabVar{D+/$+\infty$,-/$f(\alpha)$,+/$+\infty$}
\end{TAB}
\pc معامل توجيه $(T)$ هو $3$ ، إذن $f'(x_0)=3$:
\[\frac{g(x_0)}{x_0^3}=3\ \Rightarrow\ -2+4\ln x_0=0\ \Rightarrow\ \ln x_0=\frac12\ \Rightarrow\ \bx{x_0=\sqrt e}\]
تحقق: $f(\sqrt e)=3\sqrt e-2-\dfrac{2\times\frac12}{e}=3\sqrt e-2-e^{-1}$ ، ومنه $(T)$: $y=3(x-\sqrt e)+3\sqrt e-2-e^{-1}=3x-2-e^{-1}$ \ok
\QQ{4} الرسم (الوحدة $2\,cm$):
\begin{GR}{-0.3}{2.4}{-0.5}{4.6}{1.2}
\Ticks{1,2}{1,2,3,4}
\draw[cRem!70!black,thick,domain=0.2:2.3] plot(\x,{3*\x-2}) node[above,font=\scriptsize]{$(\Delta)$};
\draw[cDef,very thick,domain=0.42:2.3,samples=120] plot(\x,{3*\x-2-2*ln(\x)/(\x*\x)});
\node[cDef,font=\scriptsize] at (0.75,4.2){$(C_f)$};
\fill[cThm] (1,1) circle (1.6pt);
\end{GR}
\QQ{5} $(\star)$
\pa $H'(x)=\dfrac{\frac1x\times x-(1+\ln x)}{x^2}=\bx{-\dfrac{\ln x}{x^2}}$.
\pb على $[1;e]$: $(C_f)$ تحت $(\Delta)$ ، و $(3x-2)-f(x)=\dfrac{2\ln x}{x^2}=-2H'(x)$:
\[\mathcal A=\int_1^e\frac{2\ln x}{x^2}\,dx=-2\big[H(x)\big]_1^e=-2\left(\frac2e-1\right)=2-\frac4e\ \text{ua}\]
وحدة المساحة $4\,cm^2$ ، إذن: $\bx{\mathcal A=8-\dfrac{16}e\ cm^2}\approx2{,}11\,cm^2$.
\end{hal}

\begin{tamrin}{بكالوريا تجريبية 2026 -- تسيير واقتصاد (تلمسان) -- التمرين الرابع}
\QQ{1} الدالة $g$ معرفة على $]0;+\infty[$ بـ: $g(x)=2x^3-1+2\ln(2x)$ ، وهي متزايدة تمامًا من $-\infty$ إلى $+\infty$.
\begin{sq}
\qa بيّن أن المعادلة $g(x)=0$ تقبل حلًا وحيدًا $\alpha$ حيث $0{,}6<\alpha<0{,}7$.
\qb استنتج إشارة $g(x)$ حسب قيم $x$.
\end{sq}
\QQ{2} الدالة $f$ معرفة على $]0;+\infty[$ بـ: $f(x)=2x-1-\dfrac{\ln(2x)}{x^2}$ ، و $(C)$ تمثيلها البياني.
\begin{sq}
\qa أحسب $\Lm{0^+}f(x)$ و $\Lm{+\infty}f(x)$. (يعطى: $\Lm{+\infty}\dfrac{\ln(2x)}{x^2}=0$).
\qb بيّن أن المستقيم $(\Delta)$ ذا المعادلة $y=2x-1$ مقارب للمنحنى $(C)$.
\qc ادرس وضعية المنحنى $(C)$ بالنسبة إلى المستقيم $(\Delta)$.
\end{sq}
\QQ{3} بيّن أنه من أجل كل $x$ من $]0;+\infty[$: $f'(x)=\dfrac{g(x)}{x^3}$ ، ثم شكّل جدول تغيرات $f$.
\QQ{4} أرسم $(\Delta)$ و $(C)$.
\QQ{5} $(\star)$ الدالة $H$ معرفة على $]0;+\infty[$ بـ: $H(x)=\dfrac{1+\ln x}{x}$.
\begin{sq}
\qa بيّن أن $H$ أصلية للدالة $h:x\mapsto-\dfrac{\ln x}{x^2}$ على $]0;+\infty[$.
\qb أحسب مساحة الحيز المستوي المحدد بالمنحنى $(C)$ ، محور الفواصل والمستقيمين $x=1$ و $x=e$.
\end{sq}
\end{tamrin}

\begin{hal}
\QQ{1}
\pa $g$ مستمرة ومتزايدة تمامًا على $[0{,}6;0{,}7]$ ، و:
\[g(0{,}6)=0{,}432-1+2\ln1{,}2\approx-0{,}20<0\qquad;\qquad g(0{,}7)=0{,}686-1+2\ln1{,}4\approx0{,}36>0\]
إذن المعادلة $g(x)=0$ تقبل حلًا وحيدًا $\bx{\alpha\in]0{,}6;0{,}7[}$ ($\alpha\approx0{,}64$).
\pb $g(x)<0$ على $]0;\alpha[$ ، و $g(x)>0$ على $]\alpha;+\infty[$.
\QQ{2}
\pa
\LD{f(x)}{0^+}{+\infty} $\Car{\Lm{0^+}(2x-1)=-1\\[6pt]\Lm{0^+}\left(-\dfrac{\ln(2x)}{x^2}\right)=+\infty}$\E
لأن $\ln(2x)\to-\infty$ و $x^2\to0^+$.
\LD{f(x)}{+\infty}{+\infty} $\Car{\Lm{+\infty}(2x-1)=+\infty\\[6pt]\Lm{+\infty}\dfrac{\ln(2x)}{x^2}=0}$\E
\pb $f(x)-(2x-1)=-\dfrac{\ln(2x)}{x^2}$ و $\Lm{+\infty}\left(-\dfrac{\ln(2x)}{x^2}\right)=0$: \bxt{$(\Delta)$ مقارب مائل عند $+\infty$}.
\pc إشارة الفرق عكس إشارة $\ln(2x)$ ، و $\ln(2x)>0$ يعني $x>\dfrac12$:
\begin{FIG}
\begin{tikzpicture}[font=\normalsize]
\fill[gray!12] (0,0) rectangle (2.6,-0.9);
\fill[green!18] (2.6,0) rectangle (15,-0.9);
\fill[orange!45] (0,-0.9) rectangle (2.6,-4.6);
\fill[cyan!22] (2.6,-0.9) rectangle (15,-4.6);
\draw[thick] (0,0) rectangle (15,-4.6); \draw (0,-0.9)--(15,-0.9) (0,-1.8)--(15,-1.8) (2.6,0)--(2.6,-4.6);
\node at (1.3,-0.45){$x$};
\node at (1.3,-1.35){$-\frac{\ln(2x)}{x^2}$};
\node[align=center,font=\bfseries] at (1.3,-3.2){\ar{الوضع}\\ \ar{النسبي}};
\node at (3.3,-0.45){$0$};
\node at (8.8,-0.45){$\frac12$};
\node at (14.3,-0.45){$+\infty$};
\node at (5.3,-1.35){$+$};
\node[align=center,font=\small] at (5.3,-3.2){$(C)$\\ \ar{فوق}\\ $(\Delta)$};
\node at (12.3,-1.35){$-$};
\node[align=center,font=\small] at (12.3,-3.2){$(C)$\\ \ar{تحت}\\ $(\Delta)$};
\draw[double,thick] (3.3,-0.9)--(3.3,-4.6);
\node at (8.8,-1.35){$0$};
\draw[densely dotted] (8.8,-0.9)--(8.8,-1.13) (8.8,-1.57)--(8.8,-1.8);
\draw (7.3,-4.6) .. controls (7.375,-3.06) and (8.35,-2.136) .. (8.8,-1.8) .. controls (9.25,-2.136) and (10.225,-3.06) .. (10.3,-4.6);
\node[align=center,font=\footnotesize,inner sep=0pt] at (8.8,-3.65){\ar{$(C)$ يقطع}\\ $(\Delta)$\\ \ar{في النقطة}\\ $\left(\frac12;0\right)$};
\end{tikzpicture}
\end{FIG}
\QQ{3} $\left(\dfrac{\ln(2x)}{x^2}\right)'=\dfrac{\frac{2}{2x}x^2-2x\ln(2x)}{x^4}=\dfrac{1-2\ln(2x)}{x^3}$ ، ومنه:
\[f'(x)=2-\frac{1-2\ln(2x)}{x^3}=\frac{2x^3-1+2\ln(2x)}{x^3}=\bx{\frac{g(x)}{x^3}}\]
\begin{TAB}
\tkzTabInit[lgt=1.6,espcl=3]{$x$/1,$f'(x)$/1,$f(x)$/2}{$0$,$\alpha$,$+\infty$}
\tkzTabLine{d,-,z,+,}
\tkzTabVar{D+/$+\infty$,-/$f(\alpha)$,+/$+\infty$}
\end{TAB}
\QQ{4} الرسم:
\begin{GR}{-0.3}{2.4}{-1.3}{3.6}{1.2}
\Ticks{1,2}{-1,1,2,3}
\draw[cRem!70!black,thick,domain=0:2.3] plot(\x,{2*\x-1}) node[above,font=\scriptsize]{$(\Delta)$};
\draw[cDef,very thick,domain=0.33:2.3,samples=120] plot(\x,{2*\x-1-ln(2*\x)/(\x*\x)});
\node[cDef,font=\scriptsize] at (0.62,3.2){$(C)$};
\end{GR}
\QQ{5} $(\star)$
\pa $H'(x)=\dfrac{\frac1x\times x-(1+\ln x)}{x^2}=-\dfrac{\ln x}{x^2}=\bx{h(x)}$.
\pb على $[1;e]$: $f$ متزايدة و $f(1)=1-\ln2>0$ ، إذن $f(x)>0$. ولدينا $\dfrac{\ln(2x)}{x^2}=\dfrac{\ln2}{x^2}+\dfrac{\ln x}{x^2}$ ، ودالة أصلية لـ $f$ هي:
\[F(x)=x^2-x+\frac{\ln2}{x}+H(x)\]
\[\mathcal A=F(e)-F(1)=\left(e^2-e+\frac{\ln2}e+\frac2e\right)-\left(0+\ln2+1\right)\]
\[\mathcal A=\bx{e^2-e-1+\frac2e-\left(1-\frac1e\right)\ln2\ \text{ua}}\approx3{,}97\ \text{ua}\]
\end{hal}

\begin{tamrin}{بكالوريا تجريبية}
\textbf{I)} الدالة العددية $g$ معرفة على المجال $]0;+\infty[$ بـ: $g(x)=x\ln x-1$.
\QQ{1} شكّل جدول تغيرات الدالة $g$ على $]0;+\infty[$. (يعطى: $\Lm{0^+}x\ln x=0$).
\QQ{2}
\begin{sq}
\qa بيّن أن المعادلة $g(x)=0$ تقبل حلًا وحيدًا $\alpha$ حيث $1{,}7<\alpha<1{,}8$.
\qb استنتج حسب قيم $x$ من $]0;+\infty[$ إشارة $g(x)$.
\end{sq}
\textbf{II)} الدالة العددية $f$ معرفة على المجال $]0;+\infty[$ بـ: $f(x)=(x-1)(\ln x-1)$ ، و $(C_f)$ تمثيلها البياني.
\QQ{1}
\begin{sq}
\qa أحسب $\Lm{0^+}f(x)$ وفسّر النتيجة هندسيًا ، ثم أحسب $\Lm{+\infty}f(x)$.
\qb بيّن أنه من أجل كل $x$ من $]0;+\infty[$: $f'(x)=\dfrac{g(x)}x$.
\qc بيّن أن $f$ متناقصة تمامًا على $]0;\alpha]$ ومتزايدة تمامًا على $[\alpha;+\infty[$ ، ثم شكّل جدول تغيراتها.
\end{sq}
\QQ{2}
\begin{sq}
\qa أكتب معادلة للمماس $(T)$ للمنحنى $(C_f)$ في النقطة التي فاصلتها $1$.
\qb ادرس وضعية $(C_f)$ بالنسبة إلى $(T)$.
\end{sq}
\QQ{3}
\begin{sq}
\qa أنشئ كلًا من $(T)$ و $(C_f)$. (تعطى: $f(\alpha)\approx-0{,}33$).
\qb باستخدام $(C_f)$ ناقش بيانيًا حسب قيم الوسيط الحقيقي $m$ عدد حلول المعادلة $f(x)=-x+m$.
\end{sq}
\end{tamrin}

\begin{hal}
\textbf{I)}
\QQ{1} $g'(x)=\ln x+x\times\dfrac1x=\ln x+1$ ، و $g'(x)=0$ يعني $x=e^{-1}$ ؛ $g(e^{-1})=-e^{-1}-1$.
\LD{g(x)}{0^+}{-1} $\Car{\Lm{0^+}x\ln x=0\\[6pt]\Lm{0^+}(-1)=-1}$\E
\LD{g(x)}{+\infty}{+\infty} $\Car{\Lm{+\infty}x=+\infty\\[6pt]\Lm{+\infty}\ln x=+\infty}$\E
\begin{TAB}
\tkzTabInit[lgt=1.6,espcl=3]{$x$/1,$g'(x)$/1,$g(x)$/2}{$0$,$e^{-1}$,$+\infty$}
\tkzTabLine{d,-,z,+,}
\tkzTabVar{D+/$-1$,-/$-1-\frac1e$,+/$+\infty$}
\end{TAB}
\QQ{2}
\pa على $]0;e^{-1}]$: $g(x)\leqslant-1<0$ ، لا توجد حلول. على $[e^{-1};+\infty[$: $g$ مستمرة ومتزايدة تمامًا ، و:
\[g(1{,}7)\approx-0{,}098<0\qquad;\qquad g(1{,}8)\approx0{,}058>0\]
إذن المعادلة $g(x)=0$ تقبل حلًا وحيدًا $\bx{\alpha\in]1{,}7;1{,}8[}$.
\pb $g(x)<0$ على $]0;\alpha[$ ، و $g(x)>0$ على $]\alpha;+\infty[$.\par
\textbf{II)}
\QQ{1}
\pa
\LD{f(x)}{0^+}{+\infty} $\Car{\Lm{0^+}(x-1)=-1\\[6pt]\Lm{0^+}(\ln x-1)=-\infty}$\E
إذن محور التراتيب مقارب عمودي.
\LD{f(x)}{+\infty}{+\infty} $\Car{\Lm{+\infty}(x-1)=+\infty\\[6pt]\Lm{+\infty}(\ln x-1)=+\infty}$\E
\pb
\[f'(x)=1\times(\ln x-1)+(x-1)\times\frac1x=\ln x-\frac1x=\frac{x\ln x-1}x=\bx{\frac{g(x)}x}\]
\pc إشارة $f'$ هي إشارة $g$ ، ومنه النتيجة. و $f(\alpha)=(\alpha-1)\left(\dfrac1\alpha-1\right)=-\dfrac{(\alpha-1)^2}\alpha\approx-0{,}33$ (لأن $\ln\alpha=\dfrac1\alpha$).
\begin{TAB}
\tkzTabInit[lgt=1.6,espcl=3]{$x$/1,$f'(x)$/1,$f(x)$/2}{$0$,$\alpha$,$+\infty$}
\tkzTabLine{d,-,z,+,}
\tkzTabVar{D+/$+\infty$,-/$f(\alpha)$,+/$+\infty$}
\end{TAB}
\QQ{2}
\pa $f(1)=0$ و $f'(1)=g(1)=-1$ ، ومنه $\bx{(T):\ y=-x+1}$.
\pb $f(x)-(-x+1)=(x-1)(\ln x-1)+(x-1)=(x-1)\ln x$ ، والعددان $x-1$ و $\ln x$ لهما نفس الإشارة:
\begin{FIG}
\begin{tikzpicture}[font=\normalsize]
\fill[gray!12] (0,0) rectangle (2.6,-0.9);
\fill[green!18] (2.6,0) rectangle (15,-0.9);
\fill[orange!45] (0,-0.9) rectangle (2.6,-4.6);
\fill[cyan!22] (2.6,-0.9) rectangle (15,-4.6);
\draw[thick] (0,0) rectangle (15,-4.6); \draw (0,-0.9)--(15,-0.9) (0,-1.8)--(15,-1.8) (2.6,0)--(2.6,-4.6);
\node at (1.3,-0.45){$x$};
\node at (1.3,-1.35){$(x-1)\ln x$};
\node[align=center,font=\bfseries] at (1.3,-3.2){\ar{الوضع}\\ \ar{النسبي}};
\node at (3.3,-0.45){$0$};
\node at (8.8,-0.45){$1$};
\node at (14.3,-0.45){$+\infty$};
\node at (5.3,-1.35){$+$};
\node[align=center,font=\small] at (5.3,-3.2){$(C_f)$\\ \ar{فوق}\\ $(T)$};
\node at (12.3,-1.35){$+$};
\node[align=center,font=\small] at (12.3,-3.2){$(C_f)$\\ \ar{فوق}\\ $(T)$};
\draw[double,thick] (3.3,-0.9)--(3.3,-4.6);
\node at (8.8,-1.35){$0$};
\draw[densely dotted] (8.8,-0.9)--(8.8,-1.13) (8.8,-1.57)--(8.8,-1.8);
\draw (7.3,-4.6) .. controls (7.375,-3.06) and (8.35,-2.136) .. (8.8,-1.8) .. controls (9.25,-2.136) and (10.225,-3.06) .. (10.3,-4.6);
\node[align=center,font=\footnotesize,inner sep=0pt] at (8.8,-3.65){\ar{$(C_f)$ يمس}\\ $(T)$\\ \ar{في النقطة}\\ $(1;0)$};
\end{tikzpicture}
\end{FIG}
\QQ{3}
\pa الرسم:
\begin{GR}{-0.3}{5.6}{-1.3}{4.4}{0.85}
\Ticks{1,2,3,4,5}{-1,1,2,3,4}
\draw[cThm,thick,domain=-0.3:2.6] plot(\x,{-\x+1}) node[below,font=\scriptsize]{$(T)$};
\draw[cDef,very thick,domain=0.045:5.6,samples=150] plot(\x,{(\x-1)*(ln(\x)-1)});
\node[cDef,font=\scriptsize] at (4.6,3.6){$(C_f)$};
\end{GR}
\pb 
\[ f(x)=-x+m\iff f(x)+x=m\iff(x-1)\ln x+1=m \]
الحلول هي فواصل نقط تقاطع $(C_f)$ مع المستقيمات $y=-x+m$ الموازية لـ $(T)$ ، و $(C_f)$ فوق $(T)$ ويمسه في $(1;0)$:
\begin{itemize}
\item $m<1$: المستقيم تحت $(T)$: \bxt{لا توجد حلول}.
\item $m=1$: المستقيم هو $(T)$: \bxt{حل وحيد $x=1$}.
\item $m>1$: \bxt{حلان} ، أحدهما في $]0;1[$ والآخر في $]1;+\infty[$.
\end{itemize}
\end{hal}

% =====================================================
\section{مسائل البكالوريا -- شعبة تسيير واقتصاد}

\begin{tamrin}{بكالوريا 2026 -- الموضوع الأول -- التمرين الثالث}
يمثل الجدول التالي تغيرات الدالة $f$ المعرفة على $\R$ بـ: $f(x)=\ln\left(2x^2+ax+b\right)$ حيث $a$ و $b$ عددان حقيقيان.
\begin{TAB}
\tkzTabInit[lgt=1.6,espcl=3]{$x$/1,$f(x)$/2}{$-\infty$,$-\frac12$,$+\infty$}
\tkzTabVar{+/$+\infty$,-/$-\ln2$,+/$+\infty$}
\end{TAB}
\QQ{1}
\begin{sq}
\qa أحسب $f'(x)$ بدلالة $a$ و $b$.
\qb انطلاقًا من جدول التغيرات ، عيّن العددين $a$ و $b$.
\end{sq}
\QQ{2} نفرض أن: $f(x)=\ln\left(2x^2+2x+1\right)$.
\begin{sq}
\qa حل في $\R$ المعادلة: $f(x)=0$.
\qb حدّد إشارة كل من $f(x)$ و $f'(x)$.
\end{sq}
\QQ{3} $g$ الدالة المعرفة على $]0;+\infty[$ بـ: $g(x)=\dfrac1{f(x)}$. ادرس اتجاه تغير الدالة $g$ ثم شكّل جدول تغيراتها.
\end{tamrin}

\begin{hal}
\QQ{1}
\pa $(\ln u)'=\dfrac{u'}u$: \quad $\bx{f'(x)=\dfrac{4x+a}{2x^2+ax+b}}$.
\pb من الجدول: $f$ تقبل قيمة حدية صغرى $-\ln2$ عند $-\dfrac12$ ، إذن:
\[f'\left(-\frac12\right)=0\ \Rightarrow\ 4\times\left(-\frac12\right)+a=0\ \Rightarrow\ a=2\]
\[f\left(-\frac12\right)=-\ln2=\ln\frac12\ \Rightarrow\ 2\times\frac14+2\times\left(-\frac12\right)+b=\frac12\ \Rightarrow\ b=1\]
$\bx{a=2\ ;\ b=1}$.
\QQ{2}
\pa
\[f(x)=0\ \Rightarrow\ 2x^2+2x+1=1\ \Rightarrow\ 2x(x+1)=0\qquad;\qquad\bx{S=\{-1;0\}}\]
\pb $f(x)>0\iff2x^2+2x+1>1\iff2x(x+1)>0$ ، و $f'(x)=\dfrac{4x+2}{2x^2+2x+1}$ إشارتها إشارة $4x+2$ (المقام موجب لأن مميزه سالب):
\begin{TAB}
\tkzTabInit[lgt=1.6,espcl=2.2]{$x$/1,$f(x)$/1,$f'(x)$/1}{$-\infty$,$-1$,$-\frac12$,$0$,$+\infty$}
\tkzTabLine{,+,z,-,t,-,z,+,}
\tkzTabLine{,-,t,-,z,+,t,+,}
\end{TAB}
\QQ{3} على $]0;+\infty[$: $f(x)>0$ و $f'(x)>0$:
\[g'(x)=-\frac{f'(x)}{\big(f(x)\big)^2}<0\]
إذن \bxt{$g$ متناقصة تمامًا على $]0;+\infty[$}.
\par لدينا: $\Lm{0^+}g(x)=\bx{+\infty}$ لأنّ $\Lm{0^+}f(x)=0^+$ و $g=\dfrac1f$.\E
\par لدينا: $\Lm{+\infty}g(x)=\bx{0}$ لأنّ $\Lm{+\infty}f(x)=+\infty$.\E
\begin{TAB}
\tkzTabInit[lgt=1.6,espcl=4]{$x$/1,$g'(x)$/1,$g(x)$/2}{$0$,$+\infty$}
\tkzTabLine{d,-,}
\tkzTabVar{D+/$+\infty$,-/$0$}
\end{TAB}
\end{hal}

\begin{tamrin}{بكالوريا 2025 -- الموضوع الأول -- التمرين الرابع}
$f$ الدالة المعرفة بـ: $f(0)=0$ ، ومن أجل كل $x$ من $]0;+\infty[$: $f(x)=\dfrac14x^2(2\ln x-3)$ ، و $(C_f)$ تمثيلها البياني.
\QQ{1} أحسب $\Lm{+\infty}f(x)$.
\QQ{2}
\begin{sq}
\qa بيّن أنه من أجل كل $x$ من $]0;+\infty[$: $f'(x)=x(\ln x-1)$.
\qb ادرس اتجاه تغير الدالة $f$ ثم شكّل جدول تغيراتها.
\end{sq}
\QQ{3}
\begin{sq}
\qa بيّن أن المنحنى $(C_f)$ يقبل نقطة انعطاف $A$ ، يطلب تعيين إحداثييها.
\qb عيّن معادلة لـ $(T)$ مماس المنحنى $(C_f)$ عند النقطة $A$.
\end{sq}
\QQ{4}
\begin{sq}
\qa أحسب $f\left(e^{\frac32}\right)$ ثم ارسم كلًا من $(T)$ و $(C_f)$.
\qb ناقش بيانيًا وحسب قيم العدد الحقيقي $m$ عدد حلول المعادلة: $f(x)=m$.
\end{sq}
\QQ{5} $(\star)$ $F$ الدالة المعرفة على $]0;+\infty[$ بـ: $F(x)=\dfrac1{36}x^3(6\ln x-11)$.
\begin{sq}
\qa بيّن أنه من أجل كل $x$ من $]0;+\infty[$: $F'(x)=f(x)$.
\qb استنتج حساب $\mathcal A$ مساحة الحيز المستوي المحدد بالمنحنى $(C_f)$ والمستقيمات التي معادلاتها: $x=1$ ، $x=e$ و $y=0$.
\end{sq}
\QQ{6} $g$ الدالة المعرفة بـ: $g(0)=0$ ، ومن أجل كل عدد حقيقي غير معدوم $x$: $g(x)=\dfrac14x^2(2\ln|x|-3)$. بيّن أن الدالة $g$ زوجية ثم ارسم $(C_g)$ انطلاقًا من $(C_f)$.
\end{tamrin}

\begin{hal}
\QQ{1}
\LD{f(x)}{+\infty}{+\infty} $\Car{\Lm{+\infty}\dfrac14x^2=+\infty\\[6pt]\Lm{+\infty}(2\ln x-3)=+\infty}$\E
\QQ{2}
\pa
\[f'(x)=\frac12x(2\ln x-3)+\frac14x^2\times\frac2x=x\ln x-\frac32x+\frac12x=x\ln x-x=\bx{x(\ln x-1)}\]
\pb على $]0;+\infty[$ إشارة $f'(x)$ هي إشارة $\ln x-1$ ، وتنعدم عند $e$: $f(e)=\dfrac14e^2(2-3)=-\dfrac{e^2}4\approx-1{,}85$.
\begin{TAB}
\tkzTabInit[lgt=1.6,espcl=3]{$x$/1,$f'(x)$/1,$f(x)$/2}{$0$,$e$,$+\infty$}
\tkzTabLine{,-,z,+,}
\tkzTabVar{+/$0$,-/$-\frac{e^2}4$,+/$+\infty$}
\end{TAB}
\QQ{3}
\pa 
\[ f''(x)=(\ln x-1)+x\times\dfrac1x=\ln x \]
تنعدم عند $1$ وتغير إشارتها (سالبة على $]0;1[$ وموجبة على $]1;+\infty[$) ، و 
\[ f(1)=\dfrac14(0-3)=-\dfrac34 \]
\bxt{$A\left(1;-\frac34\right)$ نقطة انعطاف}.
\pb $f'(1)=1\times(0-1)=-1$ ، ومنه $(T)$: $y=-(x-1)-\dfrac34$ ، أي $\bx{(T):\ y=-x+\dfrac14}$.
\QQ{4}
\pa $f\left(e^{\frac32}\right)=\dfrac14e^3\left(2\times\dfrac32-3\right)=\bx{0}$: المنحنى يقطع محور الفواصل في النقطة $\left(e^{\frac32};0\right)\approx(4{,}48;0)$.
\begin{GR}{-5}{5}{-2.5}{2.6}{0.75}
\Ticks{-4,-3,-2,-1,1,2,3,4}{-2,-1,1,2}
\draw[cThm,thick,domain=-1.5:3.2] plot(\x,{-\x+0.25}) node[below,font=\scriptsize]{$(T)$};
\draw[cDef,very thick,domain=0.01:5,samples=150] plot(\x,{0.25*\x*\x*(2*ln(\x)-3)});
\draw[cRem!70!black,thick,dashed,domain=-5:-0.01,samples=150] plot(\x,{0.25*\x*\x*(2*ln(-\x)-3)});
\fill[cThm] (1,-0.75) circle (1.6pt) node[below left,font=\scriptsize]{$A$};
\node[cDef,font=\scriptsize] at (3.6,2.2){$(C_f)$}; \node[cRem!70!black,font=\scriptsize] at (-3.6,2.2){$(C_g)$};
\end{GR}
\pb الحلول هي فواصل نقط تقاطع $(C_f)$ مع المستقيم الأفقي $y=m$:
\begin{itemize}
\item $m<-\dfrac{e^2}4$: لا توجد حلول. \quad $m=-\dfrac{e^2}4$: حل وحيد $x=e$.
\item $-\dfrac{e^2}4<m<0$: حلان (أحدهما في $]0;e[$ والآخر في $\left]e;e^{\frac32}\right[$).
\item $m=0$: حلان $x=0$ و $x=e^{\frac32}$. \quad $m>0$: حل وحيد (أكبر من $e^{\frac32}$).
\end{itemize}
\QQ{5} $(\star)$
\pa
\[F'(x)=\frac1{36}\left[3x^2(6\ln x-11)+x^3\times\frac6x\right]=\frac{x^2}{36}(18\ln x-27)=\frac{x^2}4(2\ln x-3)=\bx{f(x)}\]
\pb على $[1;e]$: $2\ln x-3<0$ ، أي $f(x)<0$:
\[\mathcal A=-\big[F(e)-F(1)\big]=-\left[\frac{e^3(6-11)}{36}-\frac{-11}{36}\right]=\bx{\frac{5e^3-11}{36}\ \text{ua}}\approx2{,}48\ \text{ua}\]
\QQ{6} المجموعة $\R$ متناظرة ، و $g(-0)=g(0)$ ، ومن أجل $x\neq0$: $|-x|=|x|$ ، ومنه $g(-x)=g(x)$: \bxt{$g$ زوجية}.\par
من أجل $x\geqslant0$: $g(x)=f(x)$ ، إذن $(C_g)$ ينطبق على $(C_f)$ ، ونكمله بالتناظر بالنسبة إلى محور التراتيب (الخط المتقطع في الشكل).
\end{hal}

\begin{tamrin}{بكالوريا 2025 -- الموضوع الثاني -- التمرين الثالث}
\QQ{1} يمثل الجدول تغيرات الدالة $f$ المعرفة على $]0;+\infty[$ بـ: $f(x)=x-1-\ln x$ ($f$ متناقصة على $]0;1]$ ومتزايدة على $[1;+\infty[$ ، و $f(1)=0$).
\begin{sq}
\qa من جدول التغيرات ، حدّد إشارة $f(x)$ على $]0;+\infty[$.
\qb استنتج أنه: من أجل كل $x$ من $]0;+\infty[$: $\ln x\leqslant x-1$.
\end{sq}
\QQ{2} $(\Gamma)$ التمثيل البياني للدالة $x\mapsto\ln x$ ، و $(T)$ المماس له عند النقطة ذات الفاصلة $1$.
\begin{sq}
\qa عيّن معادلة للمماس $(T)$.
\qb استنتج مما سبق الوضع النسبي لـ $(\Gamma)$ و $(T)$.
\end{sq}
\QQ{3} $(\star)$
\begin{sq}
\qa تحقق أن الدالة $h$ المعرفة على $]0;+\infty[$ بـ: $h(x)=x\ln x-x$ أصلية للدالة $\ln$ على $]0;+\infty[$.
\qb استنتج حساب القيمة المتوسطة للدالة $\ln$ على $[1;e]$.
\end{sq}
\end{tamrin}

\begin{hal}
\QQ{1}
\pa القيمة الحدية الصغرى للدالة $f$ هي $f(1)=0$ ، إذن من أجل كل $x>0$: $\bx{f(x)\geqslant0}$ (تنعدم فقط عند $1$).
\pb $f(x)\geqslant0\iff x-1-\ln x\geqslant0\iff\bx{\ln x\leqslant x-1}$.
\QQ{2}
\pa $\ln1=0$ و $\ln'(1)=\dfrac11=1$ ، ومنه $\bx{(T):\ y=x-1}$.
\pb الفرق $\ln x-(x-1)=-f(x)\leqslant0$:
\begin{FIG}
\begin{tikzpicture}[font=\normalsize]
\fill[gray!12] (0,0) rectangle (2.6,-0.9);
\fill[green!18] (2.6,0) rectangle (15,-0.9);
\fill[orange!45] (0,-0.9) rectangle (2.6,-4.6);
\fill[cyan!22] (2.6,-0.9) rectangle (15,-4.6);
\draw[thick] (0,0) rectangle (15,-4.6); \draw (0,-0.9)--(15,-0.9) (0,-1.8)--(15,-1.8) (2.6,0)--(2.6,-4.6);
\node at (1.3,-0.45){$x$};
\node at (1.3,-1.35){$\ln x-(x-1)$};
\node[align=center,font=\bfseries] at (1.3,-3.2){\ar{الوضع}\\ \ar{النسبي}};
\node at (3.3,-0.45){$0$};
\node at (8.8,-0.45){$1$};
\node at (14.3,-0.45){$+\infty$};
\node at (5.3,-1.35){$-$};
\node[align=center,font=\small] at (5.3,-3.2){$(\Gamma)$\\ \ar{تحت}\\ $(T)$};
\node at (12.3,-1.35){$-$};
\node[align=center,font=\small] at (12.3,-3.2){$(\Gamma)$\\ \ar{تحت}\\ $(T)$};
\draw[double,thick] (3.3,-0.9)--(3.3,-4.6);
\node at (8.8,-1.35){$0$};
\draw[densely dotted] (8.8,-0.9)--(8.8,-1.13) (8.8,-1.57)--(8.8,-1.8);
\draw (7.3,-4.6) .. controls (7.375,-3.06) and (8.35,-2.136) .. (8.8,-1.8) .. controls (9.25,-2.136) and (10.225,-3.06) .. (10.3,-4.6);
\node[align=center,font=\footnotesize,inner sep=0pt] at (8.8,-3.65){\ar{$(\Gamma)$ يمس}\\ $(T)$\\ \ar{في النقطة}\\ $(1;0)$};
\end{tikzpicture}
\end{FIG}
\QQ{3} $(\star)$
\pa $h'(x)=\ln x+x\times\dfrac1x-1=\bx{\ln x}$.
\pb
\[m=\frac1{e-1}\int_1^e\ln x\,dx=\frac{h(e)-h(1)}{e-1}=\frac{(e-e)-(0-1)}{e-1}=\bx{\frac1{e-1}}\approx0{,}58\]
\end{hal}

\begin{tamrin}{بكالوريا 2024 -- الموضوع الثاني -- التمرين الرابع}
$f$ الدالة المعرفة على المجال $]0;+\infty[$ بـ: $f(x)=1+\dfrac{1+2\ln x}{x^2}$ ، و $(C_f)$ تمثيلها البياني في المعلم المتعامد والمتجانس.
\QQ{1} احسب $\Lm{0^+}f(x)$ و $\Lm{+\infty}f(x)$ ثم فسّر النتيجتين هندسيًا.
\QQ{2}
\begin{sq}
\qa بيّن أنه من أجل كل $x$ من $]0;+\infty[$: $f'(x)=\dfrac{-4\ln x}{x^3}$.
\qb ادرس اتجاه تغير الدالة $f$ ثم شكّل جدول تغيراتها.
\end{sq}
\QQ{3} بيّن أن المنحنى $(C_f)$ يقطع حامل محور الفواصل في نقطة وحيدة فاصلتها $\alpha$ حيث $0{,}52<\alpha<0{,}53$.
\QQ{4}
\begin{sq}
\qa ادرس الوضع النسبي للمنحنى $(C_f)$ والمستقيم $(\Delta)$ ذي المعادلة $y=1$.
\qb احسب $f(e)$ ثم ارسم كلًا من $(\Delta)$ و $(C_f)$.
\end{sq}
\QQ{5} $(\star)$
\begin{sq}
\qa أثبت أن $H:x\mapsto\dfrac{-3-2\ln x}x$ دالة أصلية للدالة $h:x\mapsto\dfrac{1+2\ln x}{x^2}$ على المجال $]0;+\infty[$.
\qb استنتج $\mathcal A$ مساحة الحيز المستوي المحدد بـ $(C_f)$ و $(\Delta)$ والمستقيمين اللذين معادلتاهما: $x=1$ ، $x=e$.
\end{sq}
\end{tamrin}

\begin{hal}
\QQ{1}
\LD{f(x)}{0^+}{-\infty} $\Car{\Lm{0^+}(1+2\ln x)=-\infty\\[6pt]\Lm{0^+}x^2=0^+}$\E
إذن محور التراتيب مقارب عمودي. ونكتب: $\dfrac{1+2\ln x}{x^2}=\dfrac1{x^2}+2\times\dfrac{\ln x}x\times\dfrac1x$.
\LD{f(x)}{+\infty}{1} $\Car{\Lm{+\infty}\dfrac1{x^2}=0\\[6pt]\Lm{+\infty}\dfrac{\ln x}x\times\dfrac1x=0}$\E
إذن المستقيم $y=1$ مقارب أفقي.
\QQ{2}
\pa
\[f'(x)=\frac{\frac2x\times x^2-2x(1+2\ln x)}{x^4}=\frac{2x-2x-4x\ln x}{x^4}=\bx{\frac{-4\ln x}{x^3}}\]
\pb إشارة $f'(x)$ عكس إشارة $\ln x$: \bxt{$f$ متزايدة تمامًا على $]0;1]$ ومتناقصة تمامًا على $[1;+\infty[$} ، و $f(1)=1+1=2$.
\begin{TAB}
\tkzTabInit[lgt=1.6,espcl=3]{$x$/1,$f'(x)$/1,$f(x)$/2}{$0$,$1$,$+\infty$}
\tkzTabLine{d,+,z,-,}
\tkzTabVar{D-/$-\infty$,+/$2$,-/$1$}
\end{TAB}
\QQ{3} على $[1;+\infty[$: $f(x)>1>0$ ، لا توجد حلول. على $]0;1]$: $f$ مستمرة ومتزايدة تمامًا ، و:
\[f(0{,}52)\approx-0{,}14<0\qquad;\qquad f(0{,}53)\approx0{,}04>0\]
إذن $(C_f)$ يقطع حامل محور الفواصل في نقطة وحيدة فاصلتها $\bx{\alpha\in]0{,}52;0{,}53[}$.
\QQ{4}
\pa $f(x)-1=\dfrac{1+2\ln x}{x^2}$ ، إشارتها إشارة $1+2\ln x$ ، و $1+2\ln x=0$ يعني $x=e^{-\frac12}$:
\begin{FIG}
\begin{tikzpicture}[font=\normalsize]
\fill[gray!12] (0,0) rectangle (2.6,-0.9);
\fill[green!18] (2.6,0) rectangle (15,-0.9);
\fill[orange!45] (0,-0.9) rectangle (2.6,-4.6);
\fill[cyan!22] (2.6,-0.9) rectangle (15,-4.6);
\draw[thick] (0,0) rectangle (15,-4.6); \draw (0,-0.9)--(15,-0.9) (0,-1.8)--(15,-1.8) (2.6,0)--(2.6,-4.6);
\node at (1.3,-0.45){$x$};
\node at (1.3,-1.35){$\frac{1+2\ln x}{x^2}$};
\node[align=center,font=\bfseries] at (1.3,-3.2){\ar{الوضع}\\ \ar{النسبي}};
\node at (3.3,-0.45){$0$};
\node at (8.8,-0.45){$e^{-\frac12}$};
\node at (14.3,-0.45){$+\infty$};
\node at (5.3,-1.35){$-$};
\node[align=center,font=\small] at (5.3,-3.2){$(C_f)$\\ \ar{تحت}\\ $(\Delta)$};
\node at (12.3,-1.35){$+$};
\node[align=center,font=\small] at (12.3,-3.2){$(C_f)$\\ \ar{فوق}\\ $(\Delta)$};
\draw[double,thick] (3.3,-0.9)--(3.3,-4.6);
\node at (8.8,-1.35){$0$};
\draw[densely dotted] (8.8,-0.9)--(8.8,-1.13) (8.8,-1.57)--(8.8,-1.8);
\draw (7.3,-4.6) .. controls (7.375,-3.06) and (8.35,-2.136) .. (8.8,-1.8) .. controls (9.25,-2.136) and (10.225,-3.06) .. (10.3,-4.6);
\node[align=center,font=\footnotesize,inner sep=0pt] at (8.8,-3.65){\ar{$(C_f)$ يقطع}\\ $(\Delta)$\\ \ar{في النقطة}\\ $\left(e^{-\frac12};1\right)$};
\end{tikzpicture}
\end{FIG}
\pb $f(e)=1+\dfrac{1+2}{e^2}=1+\dfrac3{e^2}\approx\bx{1{,}41}$.
\begin{GR}{-0.3}{5.6}{-1.3}{2.6}{0.9}
\Ticks{1,2,3,4,5}{-1,1,2}
\draw[cRem!70!black,thick] (-0.3,1)--(5.6,1) node[above left,font=\scriptsize]{$(\Delta)$};
\draw[cDef,very thick,domain=0.4:5.6,samples=150] plot(\x,{1+(1+2*ln(\x))/(\x*\x)});
\node[cDef,font=\scriptsize] at (1.7,2.3){$(C_f)$};
\end{GR}
\QQ{5} $(\star)$
\pa
\[H'(x)=\frac{-\frac2x\times x-(-3-2\ln x)}{x^2}=\frac{-2+3+2\ln x}{x^2}=\frac{1+2\ln x}{x^2}=\bx{h(x)}\]
\pb على $[1;e]$: $(C_f)$ فوق $(\Delta)$ ، إذن:
\[\mathcal A=\int_1^e\big(f(x)-1\big)\,dx=H(e)-H(1)=\frac{-5}e-(-3)=\bx{3-\frac5e\ \text{ua}}\approx1{,}16\ \text{ua}\]
\end{hal}

\begin{tamrin}{بكالوريا 2023 -- الموضوع الثاني -- التمرين الرابع}
$f$ الدالة المعرفة على المجال $]0;+\infty[$ بـ: $f(x)=x^2-x-\ln x$ ، و $(C_f)$ تمثيلها البياني.
\QQ{1}
\begin{sq}
\qa احسب $\Lm{0^+}f(x)$ ثم فسّر النتيجة هندسيًا.
\qb بيّن أن: $\Lm{+\infty}f(x)=+\infty$.
\end{sq}
\QQ{2}
\begin{sq}
\qa بيّن أنه من أجل كل $x$ من $]0;+\infty[$: $f'(x)=\dfrac{(2x+1)(x-1)}x$.
\qb استنتج أن الدالة $f$ متناقصة تمامًا على $]0;1]$ ومتزايدة تمامًا على $[1;+\infty[$ ثم شكّل جدول تغيراتها.
\end{sq}
\QQ{3} عيّن معادلة لـ $(T)$ المماس للمنحنى $(C_f)$ في النقطة ذات الفاصلة $2$.
\QQ{4} احسب $f(3)$ ثم ارسم $(T)$ و $(C_f)$.
\QQ{5} $(\star)$ $F$ الدالة المعرفة على $]0;+\infty[$ بـ: $F(x)=\dfrac13x^3-\dfrac12x^2+x-x\ln x$.
\begin{sq}
\qa تحقق أن $F$ أصلية للدالة $f$ على $]0;+\infty[$.
\qb احسب مساحة الحيز المستوي المحدد بالمنحنى $(C_f)$ وحامل محور الفواصل والمستقيمين اللذين معادلتاهما $x=1$ و $x=3$.
\end{sq}
\end{tamrin}

\begin{hal}
\QQ{1}
\pa
\LD{f(x)}{0^+}{+\infty} $\Car{\Lm{0^+}(x^2-x)=0\\[6pt]\Lm{0^+}(-\ln x)=+\infty}$\E
إذن محور التراتيب مقارب عمودي.
\pb حالة عدم تعيين ، نكتب: $f(x)=x^2\left(1-\dfrac1x-\dfrac{\ln x}x\times\dfrac1x\right)$.
\LD{f(x)}{+\infty}{+\infty} $\Car{\Lm{+\infty}x^2=+\infty\\[6pt]\Lm{+\infty}\left(1-\dfrac1x-\dfrac{\ln x}x\times\dfrac1x\right)=1}$\E
\QQ{2}
\pa
\[f'(x)=2x-1-\frac1x=\frac{2x^2-x-1}x=\bx{\frac{(2x+1)(x-1)}x}\]
لأن $1$ جذر لـ $2x^2-x-1$.
\pb $2x+1>0$ و $x>0$ ، إذن إشارة $f'$ هي إشارة $x-1$ ، و $f(1)=0$.
\begin{TAB}
\tkzTabInit[lgt=1.6,espcl=3]{$x$/1,$f'(x)$/1,$f(x)$/2}{$0$,$1$,$+\infty$}
\tkzTabLine{d,-,z,+,}
\tkzTabVar{D+/$+\infty$,-/$0$,+/$+\infty$}
\end{TAB}
\QQ{3} $f(2)=2-\ln2$ و $f'(2)=\dfrac{5\times1}2=\dfrac52$:
\[(T):\ y=\frac52(x-2)+2-\ln2\ \Rightarrow\ \bx{(T):\ y=\frac52x-3-\ln2}\]
\QQ{4} $f(3)=6-\ln3\approx\bx{4{,}90}$.
\begin{GR}{-0.3}{3.6}{-0.5}{5.6}{0.85}
\Ticks{1,2,3}{1,2,3,4,5}
\draw[cThm,thick,domain=1.1:3.4] plot(\x,{2.5*\x-3-ln(2)}) node[above,font=\scriptsize]{$(T)$};
\draw[cDef,very thick,domain=0.004:3.3,samples=150] plot(\x,{\x*\x-\x-ln(\x)});
\node[cDef,font=\scriptsize] at (0.6,4.3){$(C_f)$};
\fill[cThm] (2,1.307) circle (1.6pt);
\end{GR}
\QQ{5} $(\star)$
\pa $F'(x)=x^2-x+1-(\ln x+1)=x^2-x-\ln x=\bx{f(x)}$.
\pb $f(x)\geqslant0$ على $[1;3]$ ، إذن:
\[\mathcal A=F(3)-F(1)=\left(9-\frac92+3-3\ln3\right)-\left(\frac13-\frac12+1\right)=\bx{\frac{20}3-3\ln3\ \text{ua}}\approx3{,}37\ \text{ua}\]
\end{hal}

\begin{tamrin}{بكالوريا 2023 -- الموضوع الثاني -- التمرين الثالث}
\QQ{1}
\begin{sq}
\qa تحقق أنه من أجل كل عدد حقيقي $x$: $(x-2)\left(x^2-4x+3\right)=x^3-6x^2+11x-6$.
\qb حل في $\R$ المعادلة: $x^3-6x^2+11x-6=0$.
\end{sq}
\QQ{2}
\begin{sq}
\qa استنتج في المجال $]0;+\infty[$ مجموعة حلول المعادلة: $(\ln x)^3-6(\ln x)^2+11(\ln x)-6=0$.
\qb \EXP\ استنتج في $\R$ مجموعة حلول المعادلة: $e^{3x}-6e^{2x}+11e^x-6=0$.
\end{sq}
\QQ{3} حل في المجال $]2;+\infty[$ المتراجحة: $\ln\left(x^3-6x^2+11x-5\right)\geqslant0$.
\end{tamrin}

\begin{hal}
\QQ{1}
\pa $(x-2)(x^2-4x+3)=x^3-4x^2+3x-2x^2+8x-6=\bx{x^3-6x^2+11x-6}$.
\pb $x^2-4x+3=(x-1)(x-3)$ ، ومنه $P(x)=(x-1)(x-2)(x-3)=0$: $\bx{S=\{1;2;3\}}$.
\QQ{2}
\pa نضع $X=\ln x$: $X\in\{1;2;3\}$ ، أي $\bx{S=\left\{e;e^2;e^3\right\}}$.
\pb \EXP\ نضع $X=e^x$: $e^x\in\{1;2;3\}$ ، أي $\bx{S=\{0;\ln2;\ln3\}}$.
\QQ{3} $x^3-6x^2+11x-5=P(x)+1$:
\[\ln\big(P(x)+1\big)\geqslant0\ \Rightarrow\ P(x)+1\geqslant1\ \Rightarrow\ P(x)\geqslant0\]
(وعندئذ $P(x)+1>0$ فالعبارة معرفة). على $]2;+\infty[$: $x-1>0$ و $x-2>0$ ، إذن $P(x)\geqslant0$ يعني $x-3\geqslant0$:
\[\bx{S=[3;+\infty[}\]
\end{hal}

\begin{tamrin}{بكالوريا 2023 -- الموضوع الأول -- التمرين الثاني}
\QQ{1}
\begin{sq}
\qa تحقق أنه من أجل كل عدد حقيقي $x$: $(x+1)(10x-1)=10x^2+9x-1$.
\qb حل في $\R$ المعادلة: $(1-x)\left(10x^2+9x-1\right)=0$.
\end{sq}
\QQ{2}
\begin{sq}
\qa استنتج في المجال $]0;+\infty[$ مجموعة حلول المعادلة: $(1-\ln x)\left(10(\ln x)^2+9(\ln x)-1\right)=0$.
\qb \EXP\ استنتج في $\R$ مجموعة حلول المتراجحة: $\left(1-e^x\right)\left(10e^{2x}+9e^x-1\right)\geqslant0$.
\end{sq}
\QQ{3} حل في المجال $]0;+\infty[$ المتراجحة: $\ln\left(10x^2+9x\right)\geqslant0$.
\end{tamrin}

\begin{hal}
\QQ{1}
\pa $(x+1)(10x-1)=10x^2-x+10x-1=\bx{10x^2+9x-1}$.
\pb $(1-x)(x+1)(10x-1)=0$: $\bx{S=\left\{-1;\dfrac1{10};1\right\}}$.
\QQ{2}
\pa $\ln x\in\left\{-1;\dfrac1{10};1\right\}$ ، أي $\bx{S=\left\{e^{-1};e^{\frac1{10}};e\right\}}$.
\pb \EXP\ نضع $X=e^x>0$: $(1-X)(X+1)(10X-1)\geqslant0$ ، وبما أن $X+1>0$: $(1-X)(10X-1)\geqslant0$.
\begin{TAB}
\tkzTabInit[lgt=3.2,espcl=2]{$X$/1,$(1-X)(10X-1)$/1}{$0$,$\frac1{10}$,$1$,$+\infty$}
\tkzTabLine{,-,z,+,z,-,}
\end{TAB}
إذن $\dfrac1{10}\leqslant e^x\leqslant1$ ، أي $-\ln10\leqslant x\leqslant0$: \quad $\bx{S=[-\ln10;0]}$.
\QQ{3} على $]0;+\infty[$: $10x^2+9x>0$ (معرفة):
\[\ln(10x^2+9x)\geqslant0\ \Rightarrow\ 10x^2+9x\geqslant1\ \Rightarrow\ (x+1)(10x-1)\geqslant0\]
وبما أن $x+1>0$: $x\geqslant\dfrac1{10}$: \quad $\bx{S=\left[\dfrac1{10};+\infty\right[}$.
\end{hal}

\begin{tamrin}{بكالوريا 2018 -- الموضوع الأول -- التمرين الرابع}
لتكن $f$ الدالة العددية المعرفة على المجال $]-2;8[$ بـ: $f(x)=\ln(x+2)+\ln(-x+8)-\ln16$ ، و $(C_f)$ منحناها في المعلم المتعامد والمتجانس (الوحدة البيانية $2\,cm$).
\QQ{1} احسب نهايتي الدالة $f$ عند طرفي مجموعة التعريف $]-2;8[$ وفسّر النتيجتين بيانيًا.
\QQ{2} تحقّق أنه من أجل كل $x$ من $]-2;8[$: $f'(x)=\dfrac{-2x+6}{(x+2)(-x+8)}$.
\QQ{3} ادرس إشارة $f'(x)$ على المجال $]-2;8[$ وشكّل جدول تغيرات الدالة $f$.
\QQ{4} عيّن نقط تقاطع المنحنى $(C_f)$ مع محوري الإحداثيات.
\QQ{5} بيّن أنه من أجل كل $x$ من المجال $]-2;8[$: $(6-x)$ ينتمي إلى $]-2;8[$ و $f(6-x)=f(x)$ ، ثم فسّر النتيجة بيانيًا.
\QQ{6} ارسم المنحنى $(C_f)$.
\QQ{7} $(\star)$ لتكن $F$ الدالة المعرفة على $]-2;8[$ بـ: $F(x)=(x+2)\ln(x+2)+(x-8)\ln(-x+8)-2x-x\ln16$.\par
بيّن أن $F$ دالة أصلية لـ $f$ على المجال $]-2;8[$ ، ثم احسب بـ $cm^2$ مساحة الحيز المحدد بالمنحنى $(C_f)$ والمستقيمات التي معادلاتها: $y=0$ ، $x=0$ و $x=4$.
\end{tamrin}

\begin{hal}
\QQ{1}
\LD{f(x)}{-2^+}{-\infty} $\Car{\Lm{-2^+}\ln(x+2)=-\infty\\[6pt]\Lm{-2^+}\ln(-x+8)=\ln10}$\E
\LD{f(x)}{8^-}{-\infty} $\Car{\Lm{8^-}\ln(x+2)=\ln10\\[6pt]\Lm{8^-}\ln(-x+8)=-\infty}$\E
\textbf{التفسير:} المستقيمان $x=-2$ و $x=8$ مقاربان عموديان.
\QQ{2}
\[f'(x)=\frac1{x+2}+\frac{-1}{-x+8}=\frac{(-x+8)-(x+2)}{(x+2)(-x+8)}=\bx{\frac{-2x+6}{(x+2)(-x+8)}}\]
\QQ{3} المقام موجب على $]-2;8[$ ، إذن إشارة $f'$ هي إشارة $-2x+6$. و $f(3)=\ln5+\ln5-\ln16=\ln\dfrac{25}{16}\approx0{,}45$.
\begin{TAB}
\tkzTabInit[lgt=1.6,espcl=3]{$x$/1,$f'(x)$/1,$f(x)$/2}{$-2$,$3$,$8$}
\tkzTabLine{d,+,z,-,d}
\tkzTabVar{D-/$-\infty$,+/$\ln\frac{25}{16}$,-D/$-\infty$}
\end{TAB}
\QQ{4} مع محور التراتيب: $f(0)=\ln2+\ln8-\ln16=\ln\dfrac{16}{16}=0$: النقطة $\bx{O(0;0)}$.\par
مع محور الفواصل: $f(x)=0\iff\ln\big[(x+2)(8-x)\big]=\ln16$:
\[-x^2+6x+16=16\ \Rightarrow\ x(6-x)=0\qquad;\qquad\bx{(0;0)}\ \text{\ar{و}}\ \bx{(6;0)}\]
\QQ{5} $-2<x<8\iff-8<-x<2\iff-2<6-x<8$ ، و:
\[f(6-x)=\ln(6-x+2)+\ln(-6+x+8)-\ln16=\ln(8-x)+\ln(x+2)-\ln16=\bx{f(x)}\]
\textbf{التفسير:} المستقيم ذو المعادلة $x=3$ محور تناظر للمنحنى $(C_f)$.
\QQ{6} الرسم:
\begin{GR}{-2.5}{8.5}{-2.6}{1.4}{0.75}
\Ticks{-1,1,2,3,4,5,6,7}{-2,-1,1}
\draw[gray,dashed] (-2,-2.6)--(-2,1.4) (8,-2.6)--(8,1.4) (3,-2.6)--(3,1.4);
\draw[cDef,very thick,domain=-1.93:7.93,samples=200] plot(\x,{ln(\x+2)+ln(8-\x)-ln(16)});
\node[cDef,font=\scriptsize] at (3,0.85){$(C_f)$};
\end{GR}
\QQ{7} $(\star)$
\[F'(x)=\ln(x+2)+1+\ln(-x+8)+(x-8)\times\frac{-1}{-x+8}-2-\ln16\]
\[F'(x)=\ln(x+2)+\ln(-x+8)-\ln16=\bx{f(x)}\]
على $[0;4]$: $f(x)\geqslant0$ (بين الحلين $0$ و $6$) ، إذن:
\[\int_0^4f(x)\,dx=F(4)-F(0)=\big(6\ln6-4\ln4-8-4\ln16\big)-\big(2\ln2-8\ln8\big)=6\ln6-2\ln2-8\ \text{ua}\]
وحدة المساحة $4\,cm^2$: $\bx{\mathcal A=24\ln6-8\ln2-32\ cm^2}\approx5{,}46\,cm^2$.
\end{hal}

\begin{tamrin}{بكالوريا 2017 -- الموضوع الأول -- التمرين الرابع}
\textbf{I)} نعتبر الدالة $g$ المعرفة على المجال $]0;+\infty[$ كما يلي: $g(x)=x^2+3\ln x-3$.
\QQ{1} ادرس اتجاه تغيّر الدالة $g$.
\QQ{2} بيّن أن المعادلة $g(x)=0$ تقبل حلًا وحيدًا $\alpha$ حيث $1{,}40<\alpha<1{,}41$ ثم استنتج إشارة $g(x)$ حسب قيم $x$.\par
\textbf{II)} نعتبر الدالة $f$ المعرفة على المجال $]0;+\infty[$ بـ: $f(x)=x+1-\dfrac{3\ln x}x$ ، و $(C_f)$ تمثيلها البياني.
\QQ{1}
\begin{sq}
\qa احسب $\Lm{0^+}f(x)$ ، ثم فسّر النتيجة بيانيًا.
\qb احسب $\Lm{+\infty}f(x)$.
\end{sq}
\QQ{2} بيّن أنّ: من أجل كل عدد حقيقي $x$ موجب تمامًا: $f'(x)=\dfrac{g(x)}{x^2}$.
\QQ{3} استنتج اتجاه تغيّر الدالة $f$ ثم شكّل جدول تغيراتها.
\QQ{4}
\begin{sq}
\qa بيّن أنّ المستقيم $(\Delta)$ ذا المعادلة $y=x+1$ مقارب مائل للمنحنى $(C_f)$.
\qb ادرس وضعية $(C_f)$ بالنسبة إلى $(\Delta)$.
\end{sq}
\QQ{5} أنشئ المستقيم $(\Delta)$ والمنحنى $(C_f)$. (يعطى: $f(\alpha)\approx1{,}68$).
\QQ{6} $(\star)$
\begin{sq}
\qa بيّن أنّ الدالة $h$ حيث $h(x)=\dfrac12(\ln x)^2$ أصلية للدالة $x\mapsto\dfrac{\ln x}x$ على المجال $]0;+\infty[$.
\qb احسب $S$ مساحة الحيز المستوي المحدد بالمنحنى $(C_f)$ والمستقيمات التي معادلاتها: $x=1$ ، $x=e$ و $y=x+1$.
\end{sq}
\end{tamrin}

\begin{hal}
\textbf{I)}
\QQ{1} $g'(x)=2x+\dfrac3x>0$ على $]0;+\infty[$ ، إذن \bxt{$g$ متزايدة تمامًا}.
\QQ{2} $g$ مستمرة ومتزايدة تمامًا ، و:
\[g(1{,}40)\approx-0{,}031<0\qquad;\qquad g(1{,}41)\approx0{,}019>0\]
المعادلة تقبل حلًا وحيدًا $\bx{\alpha\in]1{,}40;1{,}41[}$. \quad $g(x)<0$ على $]0;\alpha[$ ، و $g(x)>0$ على $]\alpha;+\infty[$.\par
\textbf{II)}
\QQ{1}
\pa
\LD{f(x)}{0^+}{+\infty} $\Car{\Lm{0^+}(x+1)=1\\[6pt]\Lm{0^+}\left(-\dfrac{3\ln x}x\right)=+\infty}$\E
لأن $\ln x\to-\infty$ و $\dfrac1x\to+\infty$. إذن محور التراتيب مقارب عمودي.
\pb
\LD{f(x)}{+\infty}{+\infty} $\Car{\Lm{+\infty}(x+1)=+\infty\\[6pt]\Lm{+\infty}\dfrac{\ln x}x=0}$\E
\QQ{2} $\left(\dfrac{\ln x}x\right)'=\dfrac{1-\ln x}{x^2}$ ، ومنه:
\[f'(x)=1-\frac{3(1-\ln x)}{x^2}=\frac{x^2-3+3\ln x}{x^2}=\bx{\frac{g(x)}{x^2}}\]
\QQ{3} إشارة $f'$ هي إشارة $g$:
\begin{TAB}
\tkzTabInit[lgt=1.6,espcl=3]{$x$/1,$f'(x)$/1,$f(x)$/2}{$0$,$\alpha$,$+\infty$}
\tkzTabLine{d,-,z,+,}
\tkzTabVar{D+/$+\infty$,-/$f(\alpha)$,+/$+\infty$}
\end{TAB}
\QQ{4}
\pa $f(x)-(x+1)=-\dfrac{3\ln x}x$ و $\Lm{+\infty}\left(-\dfrac{3\ln x}x\right)=0$: \bxt{$(\Delta)$ مقارب مائل عند $+\infty$}.
\pb إشارة الفرق عكس إشارة $\ln x$:
\begin{FIG}
\begin{tikzpicture}[font=\normalsize]
\fill[gray!12] (0,0) rectangle (2.6,-0.9);
\fill[green!18] (2.6,0) rectangle (15,-0.9);
\fill[orange!45] (0,-0.9) rectangle (2.6,-4.6);
\fill[cyan!22] (2.6,-0.9) rectangle (15,-4.6);
\draw[thick] (0,0) rectangle (15,-4.6); \draw (0,-0.9)--(15,-0.9) (0,-1.8)--(15,-1.8) (2.6,0)--(2.6,-4.6);
\node at (1.3,-0.45){$x$};
\node at (1.3,-1.35){$-\frac{3\ln x}{x}$};
\node[align=center,font=\bfseries] at (1.3,-3.2){\ar{الوضع}\\ \ar{النسبي}};
\node at (3.3,-0.45){$0$};
\node at (8.8,-0.45){$1$};
\node at (14.3,-0.45){$+\infty$};
\node at (5.3,-1.35){$+$};
\node[align=center,font=\small] at (5.3,-3.2){$(C_f)$\\ \ar{فوق}\\ $(\Delta)$};
\node at (12.3,-1.35){$-$};
\node[align=center,font=\small] at (12.3,-3.2){$(C_f)$\\ \ar{تحت}\\ $(\Delta)$};
\draw[double,thick] (3.3,-0.9)--(3.3,-4.6);
\node at (8.8,-1.35){$0$};
\draw[densely dotted] (8.8,-0.9)--(8.8,-1.13) (8.8,-1.57)--(8.8,-1.8);
\draw (7.3,-4.6) .. controls (7.375,-3.06) and (8.35,-2.136) .. (8.8,-1.8) .. controls (9.25,-2.136) and (10.225,-3.06) .. (10.3,-4.6);
\node[align=center,font=\footnotesize,inner sep=0pt] at (8.8,-3.65){\ar{$(C_f)$ يقطع}\\ $(\Delta)$\\ \ar{في النقطة}\\ $(1;2)$};
\end{tikzpicture}
\end{FIG}
\QQ{5} الرسم:
\begin{GR}{-0.3}{5.6}{-0.3}{5.6}{0.8}
\Ticks{1,2,3,4,5}{1,2,3,4,5}
\draw[cRem!70!black,thick,domain=-0.3:4.6] plot(\x,{\x+1}) node[above,font=\scriptsize]{$(\Delta)$};
\draw[cDef,very thick,domain=0.34:5.6,samples=150] plot(\x,{\x+1-3*ln(\x)/\x});
\node[cDef,font=\scriptsize] at (4.6,3.3){$(C_f)$};
\end{GR}
\QQ{6} $(\star)$
\pa $h'(x)=\dfrac12\times2\ln x\times\dfrac1x=\bx{\dfrac{\ln x}x}$.
\pb على $[1;e]$: $(C_f)$ تحت $(\Delta)$ ، إذن:
\[S=\int_1^e3\frac{\ln x}x\,dx=3\big[h(x)\big]_1^e=3\left(\frac12-0\right)=\bx{\frac32\ \text{ua}}\]
\end{hal}

\begin{tamrin}{بكالوريا 2016 -- الموضوع الأول -- التمرين الرابع}
\textbf{I)} $g$ دالة عددية معرفة على $]0;+\infty[$ كما يلي: $g(x)=ax+b+\ln x$ حيث $a$ و $b$ عددان حقيقيان.
\QQ{1} عيّن $a$ و $b$ بحيث: $g(1)=2$ و $g'(2)=\dfrac32$.
\QQ{2} نضع: $g(x)=x+1+\ln x$.
\begin{sq}
\qa احسب $\Lm{0^+}g(x)$ و $\Lm{+\infty}g(x)$.
\qb ادرس اتجاه تغيّر الدالة $g$ ثم شكّل جدول تغيراتها.
\qc بيّن أن المعادلة $g(x)=0$ تقبل حلًا حقيقيًا وحيدًا $\alpha$ حيث: $0{,}2<\alpha<0{,}3$.
\qd حدّد تبعًا لقيم العدد الحقيقي $x$ إشارة $g(x)$ على المجال $]0;+\infty[$.
\end{sq}
\textbf{II)} نعتبر الدالة $f$ المعرّفة على المجال $]0;+\infty[$ كما يلي: $f(x)=\dfrac{x\ln x}{x+1}$.
\QQ{1} بيّن أنه من أجل كل $x$ من المجال $]0;+\infty[$: $f'(x)=\dfrac{g(x)}{(x+1)^2}$ ثم استنتج اتجاه تغيّر الدالة $f$.
\QQ{2} احسب $\Lm{0^+}f(x)$ و $\Lm{+\infty}f(x)$. (يُعطى: $\Lm{0^+}x\ln x=0$).
\QQ{3} تحقق أن: $f(\alpha)=-\alpha$ ثم شكّل جدول تغيرات الدالة $f$.
\QQ{4} احسب $f(1)$ و $f(5)$ ثم ارسم $(C_f)$ على المجال $]0;5]$.
\end{tamrin}

\begin{hal}
\textbf{I)}
\QQ{1} $g'(x)=a+\dfrac1x$: \quad $g(1)=a+b=2$ و $g'(2)=a+\dfrac12=\dfrac32$ ، ومنه $\bx{a=1\ ;\ b=1}$.
\QQ{2}
\pa
\LD{g(x)}{0^+}{-\infty} $\Car{\Lm{0^+}(x+1)=1\\[6pt]\Lm{0^+}\ln x=-\infty}$\E
\LD{g(x)}{+\infty}{+\infty} $\Car{\Lm{+\infty}(x+1)=+\infty\\[6pt]\Lm{+\infty}\ln x=+\infty}$\E
\pb $g'(x)=1+\dfrac1x>0$: \bxt{$g$ متزايدة تمامًا على $]0;+\infty[$}.
\begin{TAB}
\tkzTabInit[lgt=1.6,espcl=4]{$x$/1,$g'(x)$/1,$g(x)$/2}{$0$,$+\infty$}
\tkzTabLine{d,+,}
\tkzTabVar{D-/$-\infty$,+/$+\infty$}
\end{TAB}
\pc $g(0{,}2)=1{,}2+\ln0{,}2\approx-0{,}41<0$ و $g(0{,}3)=1{,}3+\ln0{,}3\approx0{,}10>0$: حل وحيد $\bx{\alpha\in]0{,}2;0{,}3[}$.
\pd $g(x)<0$ على $]0;\alpha[$ ، و $g(x)>0$ على $]\alpha;+\infty[$.\par
\textbf{II)}
\QQ{1}
\[f'(x)=\frac{(\ln x+1)(x+1)-x\ln x}{(x+1)^2}=\frac{x\ln x+\ln x+x+1-x\ln x}{(x+1)^2}=\frac{x+1+\ln x}{(x+1)^2}=\bx{\frac{g(x)}{(x+1)^2}}\]
إذن \bxt{$f$ متناقصة تمامًا على $]0;\alpha]$ ومتزايدة تمامًا على $[\alpha;+\infty[$}.
\QQ{2}
\LD{f(x)}{0^+}{0} $\Car{\Lm{0^+}x\ln x=0\\[6pt]\Lm{0^+}(x+1)=1}$\E
نكتب: $f(x)=\ln x\times\dfrac x{x+1}$.
\LD{f(x)}{+\infty}{+\infty} $\Car{\Lm{+\infty}\ln x=+\infty\\[6pt]\Lm{+\infty}\dfrac x{x+1}=1}$\E
\QQ{3} $g(\alpha)=0$ معناه $\ln\alpha=-\alpha-1$ ، ومنه:
\[f(\alpha)=\frac{\alpha(-\alpha-1)}{\alpha+1}=\frac{-\alpha(\alpha+1)}{\alpha+1}=\bx{-\alpha}\]
\begin{TAB}
\tkzTabInit[lgt=1.6,espcl=3]{$x$/1,$f'(x)$/1,$f(x)$/2}{$0$,$\alpha$,$+\infty$}
\tkzTabLine{d,-,z,+,}
\tkzTabVar{D+/$0$,-/$-\alpha$,+/$+\infty$}
\end{TAB}
\QQ{4} $f(1)=\bx{0}$ ؛ \quad $f(5)=\dfrac{5\ln5}6\approx\bx{1{,}34}$.
\begin{GR}{-0.3}{5.6}{-0.6}{1.8}{1}
\Ticks{1,2,3,4,5}{1}
\draw[cDef,very thick,domain=0.01:5,samples=150] plot(\x,{\x*ln(\x)/(\x+1)});
\fill[cDef] (5,1.341) circle (1.6pt);
\node[cDef,font=\scriptsize] at (3.5,1.25){$(C_f)$};
\end{GR}
\end{hal}

\begin{tamrin}{بكالوريا 2016 -- الموضوع الثاني -- التمرين الرابع}
\textbf{I)} نعتبر الدالة $g$ المعرفة على المجال $]0;+\infty[$ بـ: $g(x)=-4+2x(1+\ln x)$.
\QQ{1} احسب $\Lm{0^+}g(x)$ و $\Lm{+\infty}g(x)$. (تُعطى: $\Lm{0^+}x\ln x=0$).
\QQ{2} ادرس اتجاه تغيّر الدالة $g$ على $]0;+\infty[$ ثم شكّل جدول تغيراتها.
\QQ{3} بيّن أن المعادلة $g(x)=0$ تقبل حلًا وحيدًا $\alpha$ حيث: $1{,}4<\alpha<1{,}5$.
\QQ{4} حدّد إشارة $g(x)$ على المجال $]0;+\infty[$.\par
\textbf{II)} نعتبر الدالة $f$ المعرفة على المجال $]0;+\infty[$ بـ: $f(x)=(2x-4)\ln x$ ، و $(C_f)$ تمثيلها البياني.
\QQ{1}
\begin{sq}
\qa احسب $\Lm{0^+}f(x)$. فسّر النتيجة هندسيًا.
\qb احسب $\Lm{+\infty}f(x)$.
\end{sq}
\QQ{2}
\begin{sq}
\qa بيّن أنه من أجل كل $x$ من المجال $]0;+\infty[$: $f'(x)=\dfrac{g(x)}x$.
\qb استنتج اتجاه تغيّر الدالة $f$ ثم شكّل جدول تغيراتها.
\end{sq}
\QQ{3} عيّن نقط تقاطع $(C_f)$ مع حامل محور الفواصل.
\QQ{4}
\begin{sq}
\qa اكتب معادلة للمماس $(T)$ للمنحنى $(C_f)$ عند النقطة ذات الفاصلة $1$.
\qb أنشئ $(T)$ و $(C_f)$. (تُعطى: $f(\alpha)\approx-0{,}41$).
\end{sq}
\QQ{5} $(\star)$ نعتبر الدالة $F$ المعرفة على $]0;+\infty[$ بـ: $F(x)=\left(x^2-4x\right)\ln x-\dfrac12x^2+4x$.
\begin{sq}
\qa بيّن أن $F$ دالة أصلية للدالة $f$ على المجال $]0;+\infty[$.
\qb احسب مساحة الحيز المستوي المحدد بالمنحنى $(C_f)$ والمستقيمات التي معادلاتها: $y=0$ ، $x=1$ و $x=2$.
\end{sq}
\end{tamrin}

\begin{hal}
\textbf{I)}
\QQ{1} $g(x)=-4+2x+2x\ln x$.
\LD{g(x)}{0^+}{-4} $\Car{\Lm{0^+}(-4+2x)=-4\\[6pt]\Lm{0^+}2x\ln x=0}$\E
\LD{g(x)}{+\infty}{+\infty} $\Car{\Lm{+\infty}2x=+\infty\\[6pt]\Lm{+\infty}(1+\ln x)=+\infty}$\E
\QQ{2} $g'(x)=2(1+\ln x)+2x\times\dfrac1x=2\ln x+4$ ، و 
\[ g'(x)=0\iff\ln x=-2\iff x=e^{-2} \]
\quad $g(e^{-2})=-4+2e^{-2}(1-2)=-4-2e^{-2}$.
\begin{TAB}
\tkzTabInit[lgt=1.6,espcl=3]{$x$/1,$g'(x)$/1,$g(x)$/2}{$0$,$e^{-2}$,$+\infty$}
\tkzTabLine{d,-,z,+,}
\tkzTabVar{D+/$-4$,-/$-4-2e^{-2}$,+/$+\infty$}
\end{TAB}
\QQ{3} على $]0;e^{-2}]$: $g(x)<-4<0$. على $[e^{-2};+\infty[$: $g$ مستمرة ومتزايدة تمامًا ، و:
\[g(1{,}4)\approx-0{,}26<0\qquad;\qquad g(1{,}5)\approx0{,}22>0\]
حل وحيد $\bx{\alpha\in]1{,}4;1{,}5[}$.
\QQ{4} $g(x)<0$ على $]0;\alpha[$ ، و $g(x)>0$ على $]\alpha;+\infty[$.\par
\textbf{II)}
\QQ{1}
\pa
\LD{f(x)}{0^+}{+\infty} $\Car{\Lm{0^+}(2x-4)=-4\\[6pt]\Lm{0^+}\ln x=-\infty}$\E
إذن محور التراتيب مقارب عمودي.
\pb
\LD{f(x)}{+\infty}{+\infty} $\Car{\Lm{+\infty}(2x-4)=+\infty\\[6pt]\Lm{+\infty}\ln x=+\infty}$\E
\QQ{2}
\pa
\[f'(x)=2\ln x+(2x-4)\times\frac1x=\frac{2x\ln x+2x-4}x=\bx{\frac{g(x)}x}\]
\pb إشارة $f'$ هي إشارة $g$: \bxt{$f$ متناقصة تمامًا على $]0;\alpha]$ ومتزايدة تمامًا على $[\alpha;+\infty[$}.
\begin{TAB}
\tkzTabInit[lgt=1.6,espcl=3]{$x$/1,$f'(x)$/1,$f(x)$/2}{$0$,$\alpha$,$+\infty$}
\tkzTabLine{d,-,z,+,}
\tkzTabVar{D+/$+\infty$,-/$f(\alpha)$,+/$+\infty$}
\end{TAB}
\QQ{3} $f(x)=0\iff2x-4=0$ أو $\ln x=0$ ، أي $x=2$ أو $x=1$: النقطتان $\bx{(1;0)}$ و $\bx{(2;0)}$.
\QQ{4}
\pa $f(1)=0$ و $f'(1)=g(1)=-2$: \quad $\bx{(T):\ y=-2x+2}$.
\pb الرسم:
\begin{GR}{-0.3}{3.6}{-1.3}{3.6}{0.95}
\Ticks{1,2,3}{-1,1,2,3}
\draw[cThm,thick,domain=-0.2:2.2] plot(\x,{-2*\x+2}) node[below,font=\scriptsize]{$(T)$};
\draw[cDef,very thick,domain=0.19:3.5,samples=150] plot(\x,{(2*\x-4)*ln(\x)});
\node[cDef,font=\scriptsize] at (3,3.1){$(C_f)$};
\end{GR}
\QQ{5} $(\star)$
\pa
\[F'(x)=(2x-4)\ln x+\frac{x^2-4x}x-x+4=(2x-4)\ln x+x-4-x+4=\bx{f(x)}\]
\pb على $[1;2]$: $2x-4\leqslant0$ و $\ln x\geqslant0$ ، إذن $f(x)\leqslant0$:
\[\mathcal A=-\big[F(2)-F(1)\big]=-\left[(-4\ln2-2+8)-\left(-\frac12+4\right)\right]=\bx{4\ln2-\frac52\ \text{ua}}\approx0{,}27\ \text{ua}\]
\end{hal}

% =====================================================
\section{أسئلة الاختيار من متعدد (من البكالوريا والاختبارات)}

\begin{tamrin}{عيّن الاقتراح الصحيح الوحيد من بين الاقتراحات الثلاثة مع التبرير}
\QQ{1} (بكالوريا 2026) مجموعة حلول المعادلة $(\ln x-1)(\ln x-2)=0$ في $]0;+\infty[$ هي:\ABC{$\{1;2\}$}{$\{2;e\}$}{$\{e;e^2\}$}
\QQ{2} (بكالوريا 2023) $f(x)=1+\ln(x+1)$ على $]-1;+\infty[$. معادلة مماس $(C_f)$ عند النقطة ذات الفاصلة $0$ هي:\ABC{$y=x+1$}{$y=x$}{$y=x-1$}
\QQ{3} $(\star)$ (بكالوريا 2023) $g(x)=2x-\dfrac1x$ على $]0;+\infty[$ ، دالتها الأصلية $G$ التي تنعدم من أجل $1$ هي:\ABC{$x^2+1-\ln x$}{$-x^2+1-\ln x$}{$x^2-1-\ln x$}
\QQ{4} $(\star)$ \EXP\ (بكالوريا 2023) الدالة الأصلية على $\R$ للدالة $f(x)=\dfrac{2e^x}{e^x+1}$ والتي تنعدم من أجل $0$ هي:\ABC{$-2\ln(e^x+1)+\ln4$}{$2\ln(e^x+1)-\ln4$}{$\ln\left(\dfrac{e^x+1}2\right)$}
\QQ{5} \EXP\ (بكالوريا 2016) $f(x)=e^x+2e^{-x}-3$. حلا المعادلة $f(x)=0$ هما:\ABC{$0$ و $\ln2$}{$0$ و $-\ln2$}{$0$ و $\ln3$}
\QQ{6} (بكالوريا تجريبية 2026) $(u_n)$ متتالية حسابية أساسها $r=\ln2$ و $u_1=\ln4$. عبارة $u_n$ بدلالة $n$ هي:\ABC{$\ln4+n\ln2$}{$(n+1)\ln2$}{$\ln4-n\ln2$}
\QQ{7} $(\star)$ \EXP\ (بكالوريا تجريبية 2026) القيمة المتوسطة $m$ للدالة $f(x)=\dfrac{\ln(2)e^x}{e^x+1}$ على $[0;\ln2]$ هي:\ABC{$\ln\left(\dfrac32\right)$}{$-\ln\left(\dfrac32\right)$}{$\ln\left(\dfrac23\right)$}
\QQ{8} (بكالوريا تجريبية 2026) $\Lm{+\infty}\dfrac{\ln(3x)}{1+3\ln x}$ هي:\ABC{$0$}{$1$}{$\dfrac13$}
\QQ{9} (بكالوريا تجريبية 2025) معادلة المماس لمنحنى الدالة $f(x)=x^2-1-\ln x$ عند النقطة $A(1;0)$ هي:\ABC{$y=2x-2$}{$y=x-1$}{$y=x+2$}
\QQ{10} (بكالوريا تجريبية 2025) مجموعة حلول المتراجحة $2\ln x-1>1$ في $\R$ هي:\ABC{$\left]\dfrac12;+\infty\right[$}{$]1;+\infty[$}{$]e;+\infty[$}
\QQ{11} $(\star)$ (بكالوريا تجريبية 2025) قيمة العدد $\displaystyle\int_1^e\left(\frac{x+1}x\right)dx$ هي:\ABC{$e$}{$0$}{$e-1$}
\QQ{12} (بكالوريا تجريبية 2025 -- صحيح أم خاطئ) مجموعة حلول المعادلة $\ln(1-x)=3$ في المجال $]-\infty;1[$ هي $\left\{e^3-1\right\}$.
\end{tamrin}

\begin{hal}
\QQ{1} $\ln x=1$ أو $\ln x=2$ ، أي $x=e$ أو $x=e^2$. \quad \bxt{الجواب (جـ)}.
\QQ{2} $f(0)=1$ ، $f'(x)=\dfrac1{x+1}$ ومنه $f'(0)=1$: $y=x+1$. \quad \bxt{الجواب (أ)}.
\QQ{3} $G(x)=x^2-\ln x+c$ و $G(1)=1+c=0$ ، ومنه $c=-1$: $G(x)=x^2-1-\ln x$. \quad \bxt{الجواب (جـ)}.
\QQ{4} \EXP\ $F(x)=2\ln(e^x+1)+c$ ، $F(0)=2\ln2+c=0$ ومنه $c=-2\ln2=-\ln4$. \quad \bxt{الجواب (ب)}.
\QQ{5} \EXP\ نضع $X=e^x$: $X+\dfrac2X-3=0\iff X^2-3X+2=0\iff X=1$ أو $X=2$ ، أي $x=0$ أو $x=\ln2$. \quad \bxt{الجواب (أ)}.
\QQ{6} $u_n=u_1+(n-1)r=\ln4+(n-1)\ln2=2\ln2+(n-1)\ln2=(n+1)\ln2$. \quad \bxt{الجواب (ب)}.
\QQ{7} \EXP\ $m=\dfrac1{\ln2}\displaystyle\int_0^{\ln2}\dfrac{\ln2\,e^x}{e^x+1}\,dx=\Big[\ln(e^x+1)\Big]_0^{\ln2}=\ln3-\ln2=\ln\dfrac32$. \quad \bxt{الجواب (أ)}.
\QQ{8} نكتب: $\dfrac{\ln3+\ln x}{1+3\ln x}=\dfrac{\dfrac{\ln3}{\ln x}+1}{\dfrac1{\ln x}+3}$.
\par لدينا: $\Lm{+\infty}\dfrac{\ln(3x)}{1+3\ln x}=\bx{\dfrac13}$ لأنّ: \quad $\Car{\Lm{+\infty}\left(\dfrac{\ln3}{\ln x}+1\right)=1\\[6pt]\Lm{+\infty}\left(\dfrac1{\ln x}+3\right)=3}$\E
\bxt{الجواب (جـ)}.
\QQ{9} $f(1)=0$ ، $f'(x)=2x-\dfrac1x$ ومنه $f'(1)=1$: $y=x-1$. \quad \bxt{الجواب (ب)}.
\QQ{10} $D=]0;+\infty[$ ؛ $2\ln x>2\iff\ln x>1\iff x>e$. \quad \bxt{الجواب (جـ)}.
\QQ{11} $\displaystyle\int_1^e\left(1+\frac1x\right)dx=\big[x+\ln x\big]_1^e=(e+1)-(1+0)=e$. \quad \bxt{الجواب (أ)}.
\QQ{12} $\ln(1-x)=3\iff1-x=e^3\iff x=1-e^3$ (و $1-e^3<1$). مجموعة الحلول $\left\{1-e^3\right\}$: \bxt{خاطئ}.
\end{hal}

\begin{BOX}{cSec}{}
\centering{\Large\bfseries بالتوفيق في شهادة البكالوريا}
\end{BOX}
\end{document}
